% !TEX program = xelatex
\documentclass[12pt,a4paper]{ctexbook}

% ─── 页面与字体 ───
\usepackage[top=2.5cm,bottom=2.5cm,left=2.5cm,right=2.5cm]{geometry}
\usepackage{fontspec}
\usepackage{iffont}
\IfFontExistsTF{Consolas}
  {\setmonofont{Consolas}}
  {\setmonofont{DejaVu Sans Mono}[Scale=MatchLowercase]}

% ─── 数学 ───
\usepackage{amsmath,amssymb}

% ─── 颜色 ───
\usepackage[dvipsnames,svgnames,x11names]{xcolor}
\definecolor{codebg}{HTML}{F5F5F5}
\definecolor{codeframe}{HTML}{E0E0E0}
\definecolor{cppkeyword}{HTML}{1A1AFF}
\definecolor{cppcomment}{HTML}{2E8B2E}
\definecolor{cppstring}{HTML}{B22222}
\definecolor{linenumgray}{HTML}{999999}
\definecolor{algheadbg}{HTML}{1A3C5E}
\definecolor{luaheadbg}{HTML}{00007D}
\definecolor{pythonheadbg}{HTML}{306998}
\definecolor{angelscriptbg}{HTML}{8B0000}
\definecolor{chaibg}{HTML}{E65100}
\definecolor{squirrelbg}{HTML}{6A1B9A}
\definecolor{csharpheadbg}{HTML}{68217A}
\definecolor{warnheadbg}{HTML}{B71C1C}
\definecolor{solheadbg}{HTML}{0D47A1}
\definecolor{blueprintheadbg}{HTML}{0E4DA4}
\definecolor{verseheadbg}{HTML}{1B5E20}
\definecolor{tclheadbg}{HTML}{D84315}
\definecolor{jsheadbg}{HTML}{F9A825}
\definecolor{guileheadbg}{HTML}{4E342E}
\definecolor{rhaiheadbg}{HTML}{E53935}
\definecolor{vbaheadbg}{HTML}{0072C6}

% ─── 代码高亮 ───
\usepackage{listings}

% C++ 风格（默认）
\lstdefinestyle{cppstyle}{
  language=C++,
  basicstyle=\small\ttfamily,
  keywordstyle=\color{cppkeyword}\bfseries,
  commentstyle=\color{cppcomment}\itshape,
  stringstyle=\color{cppstring},
  frame=single,
  rulecolor=\color{codeframe},
  framesep=6pt,
  xleftmargin=1em,
  xrightmargin=1em,
  numbers=left,
  numbersep=8pt,
  numberstyle=\tiny\color{linenumgray},
  backgroundcolor=\color{codebg},
  breaklines=true,
  tabsize=4,
  columns=flexible,
  keepspaces=true,
  escapeinside={(*@}{@*)},
  showstringspaces=false,
  emph={std,string,vector,map,unordered\_map,size\_t,int64\_t,uint32\_t,shared\_ptr,unique\_ptr,function,optional,variant,tuple,pair,array,span,string\_view},
  emphstyle=\color{teal!80!black},
  morekeywords={constexpr,consteval,constinit,co\_await,co\_yield,co\_return,requires,concept,noexcept,override,final,nullptr,nullptr\_t,static\_assert,thread\_local,mutable,explicit,decltype,auto},
}

% Lua 风格
\lstdefinestyle{luastyle}{
  language=[5.3]Lua,
  basicstyle=\small\ttfamily,
  keywordstyle=\color{luaheadbg}\bfseries,
  commentstyle=\color{cppcomment}\itshape,
  stringstyle=\color{cppstring},
  frame=single,
  rulecolor=\color{codeframe},
  framesep=6pt,
  xleftmargin=1em,
  xrightmargin=1em,
  numbers=left,
  numbersep=8pt,
  numberstyle=\tiny\color{linenumgray},
  backgroundcolor=\color{codebg},
  breaklines=true,
  tabsize=4,
  columns=flexible,
  keepspaces=true,
  escapeinside={(*@}{@*)},
  showstringspaces=false,
  morekeywords={local,function,end,if,then,else,elseif,for,while,do,repeat,until,return,in,nil,true,false,and,or,not,break,goto},
}

% Python 风格
\lstdefinestyle{pythonstyle}{
  language=Python,
  basicstyle=\small\ttfamily,
  keywordstyle=\color{pythonheadbg}\bfseries,
  commentstyle=\color{cppcomment}\itshape,
  stringstyle=\color{cppstring},
  frame=single,
  rulecolor=\color{codeframe},
  framesep=6pt,
  xleftmargin=1em,
  xrightmargin=1em,
  numbers=left,
  numbersep=8pt,
  numberstyle=\tiny\color{linenumgray},
  backgroundcolor=\color{codebg},
  breaklines=true,
  tabsize=4,
  columns=flexible,
  keepspaces=true,
  escapeinside={(*@}{@*)},
  showstringspaces=false,
  emph={self,cls,None,True,False},
  emphstyle=\color{teal!80!black},
}

% C# 风格
\lstdefinestyle{csharpstyle}{
  language=[Sharp]C,
  basicstyle=\small\ttfamily,
  keywordstyle=\color{csharpheadbg}\bfseries,
  commentstyle=\color{cppcomment}\itshape,
  stringstyle=\color{cppstring},
  frame=single,
  rulecolor=\color{codeframe},
  framesep=6pt,
  xleftmargin=1em,
  xrightmargin=1em,
  numbers=left,
  numbersep=8pt,
  numberstyle=\tiny\color{linenumgray},
  backgroundcolor=\color{codebg},
  breaklines=true,
  tabsize=4,
  columns=flexible,
  keepspaces=true,
  escapeinside={(*@}{@*)},
  showstringspaces=false,
  emph={string,int,bool,var,void,object,dynamic,List,Dictionary,Task},
  emphstyle=\color{teal!80!black},
  morekeywords={async,await,yield,namespace,using,class,interface,struct,enum,delegate,event,get,set,init,record,required,partial,override,virtual,abstract,sealed,static,readonly,const,ref,out,in,params,new,typeof,is,as,switch,when,foreach,lock},
}

\lstset{style=cppstyle}

% ─── 表格 ───
\usepackage{booktabs}
\usepackage{tabularx}
\usepackage{multirow}
\usepackage{array}
\usepackage{longtable}

% ─── 长文本拆分 ───
\usepackage{seqsplit}

% ─── 彩色文本框 ───
\usepackage[most]{tcolorbox}

\newenvironment{guideline}[1][]{%
  \begin{tcolorbox}[enhanced,breakable,
    colback=blue!5!white,colframe=blue!75!black,
    boxed title style={colback=blue!75!black},
    attach boxed title to top left={yshift=-2mm,xshift=2mm},
    arc=2mm,boxrule=0.8mm,
    title={脚本集成要点},#1]}%
  {\end{tcolorbox}}

\newenvironment{compare}[1][]{%
  \begin{tcolorbox}[enhanced,breakable,
    colback=green!3!white,colframe=green!50!black,
    boxed title style={colback=green!50!black},
    attach boxed title to top left={yshift=-2mm,xshift=2mm},
    arc=2mm,boxrule=0.8mm,
    title={方案对比},#1]}%
  {\end{tcolorbox}}

\newenvironment{warning}[1][]{%
  \begin{tcolorbox}[enhanced,breakable,
    colback=red!5!white,colframe=red!75!black,
    boxed title style={colback=red!75!black},
    attach boxed title to top left={yshift=-2mm,xshift=2mm},
    arc=2mm,boxrule=0.8mm,
    title={常见陷阱},#1]}%
  {\end{tcolorbox}}

\newenvironment{note}[1][]{%
  \begin{tcolorbox}[enhanced,breakable,
    colback=orange!5!white,colframe=orange!80!black,
    boxed title style={colback=orange!80!black},
    attach boxed title to top left={yshift=-2mm,xshift=2mm},
    arc=2mm,boxrule=0.8mm,
    title={注意},#1]}%
  {\end{tcolorbox}}

\newenvironment{bestpractice}[1][]{%
  \begin{tcolorbox}[enhanced,breakable,
    colback=purple!5!white,colframe=purple!60!black,
    boxed title style={colback=purple!60!black},
    attach boxed title to top left={yshift=-2mm,xshift=2mm},
    arc=2mm,boxrule=0.8mm,
    title={最佳实践},#1]}%
  {\end{tcolorbox}}

% ─── 标题格式 ───
\usepackage{titlesec}
\titleformat{\chapter}[display]
  {\normalfont\huge\bfseries\color{algheadbg}}
  {\chaptertitlename\ \thechapter}{20pt}{\Huge}
\titleformat{\section}
  {\normalfont\Large\bfseries\color{blue!70!black}}
  {\thesection}{1em}{}
\titleformat{\subsection}
  {\normalfont\large\bfseries\color{blue!60!black}}
  {\thesubsection}{1em}{}

% ─── 页眉页脚 ───
\usepackage{fancyhdr}
\pagestyle{fancy}
\fancyhf{}
\fancyhead[LE,RO]{\thepage}
\fancyhead[RE]{\leftmark}
\fancyhead[LO]{\rightmark}

% ─── 列表 ───
\usepackage{enumitem}
\setlist[itemize]{leftmargin=2em,itemsep=2pt}
\setlist[enumerate]{leftmargin=2em,itemsep=2pt}

% ─── 超链接 ───
\usepackage{hyperref}
\hypersetup{
  colorlinks=true,
  linkcolor=blue!70!black,
  urlcolor=blue!60!black,
  pdfauthor={C/C++ Scripting Integration},
  pdftitle={C/C++ scripting integration comparison reference manual},
}

% ─── 自定义命令 ───
\newcommand{\cpp}[1]{C++#1}
\newcommand{\Fn}[1]{\texttt{#1()}}
\newcommand{\Tp}[1]{\texttt{#1}}

% ─── 语言标签 ───
\newcommand{\cpplabel}{%
  \noindent\begin{tcolorbox}[colback=algheadbg,coltext=white,
    sharp corners,boxrule=0pt,width=3.0em,height=1.4em,
    halign=center,valign=center,nobeforeafter,left=0pt,right=0pt,
    top=0pt,bottom=0pt]\small\bfseries C++\end{tcolorbox}}

\newcommand{\lualabel}{%
  \noindent\begin{tcolorbox}[colback=luaheadbg,coltext=white,
    sharp corners,boxrule=0pt,width=3.0em,height=1.4em,
    halign=center,valign=center,nobeforeafter,left=0pt,right=0pt,
    top=0pt,bottom=0pt]\small\bfseries Lua\end{tcolorbox}}

\newcommand{\pythonlabel}{%
  \noindent\begin{tcolorbox}[colback=pythonheadbg,coltext=white,
    sharp corners,boxrule=0pt,width=4.5em,height=1.4em,
    halign=center,valign=center,nobeforeafter,left=0pt,right=0pt,
    top=0pt,bottom=0pt]\small\bfseries Python\end{tcolorbox}}

\newcommand{\angelscriptlabel}{%
  \noindent\begin{tcolorbox}[colback=angelscriptbg,coltext=white,
    sharp corners,boxrule=0pt,width=6.5em,height=1.4em,
    halign=center,valign=center,nobeforeafter,left=0pt,right=0pt,
    top=0pt,bottom=0pt]\small\bfseries AngelScript\end{tcolorbox}}

\newcommand{\chaiscriptlabel}{%
  \noindent\begin{tcolorbox}[colback=chaibg,coltext=white,
    sharp corners,boxrule=0pt,width=6.0em,height=1.4em,
    halign=center,valign=center,nobeforeafter,left=0pt,right=0pt,
    top=0pt,bottom=0pt]\small\bfseries ChaiScript\end{tcolorbox}}

\newcommand{\squirrellabel}{%
  \noindent\begin{tcolorbox}[colback=squirrelbg,coltext=white,
    sharp corners,boxrule=0pt,width=5.0em,height=1.4em,
    halign=center,valign=center,nobeforeafter,left=0pt,right=0pt,
    top=0pt,bottom=0pt]\small\bfseries Squirrel\end{tcolorbox}}

\newcommand{\csharplabel}{%
  \noindent\begin{tcolorbox}[colback=csharpheadbg,coltext=white,
    sharp corners,boxrule=0pt,width=3.2em,height=1.4em,
    halign=center,valign=center,nobeforeafter,left=0pt,right=0pt,
    top=0pt,bottom=0pt]\small\bfseries C\#\end{tcolorbox}}

\newcommand{\sollabel}{%
  \noindent\begin{tcolorbox}[colback=solheadbg,coltext=white,
    sharp corners,boxrule=0pt,width=4.5em,height=1.4em,
    halign=center,valign=center,nobeforeafter,left=0pt,right=0pt,
    top=0pt,bottom=0pt]\small\bfseries sol2\end{tcolorbox}}

\newcommand{\pybindlabel}{%
  \noindent\begin{tcolorbox}[colback=pythonheadbg!80!black,coltext=white,
    sharp corners,boxrule=0pt,width=5.0em,height=1.4em,
    halign=center,valign=center,nobeforeafter,left=0pt,right=0pt,
    top=0pt,bottom=0pt]\small\bfseries pybind11\end{tcolorbox}}

\newcommand{\blueprintlabel}{%
  \noindent\begin{tcolorbox}[colback=blueprintheadbg,coltext=white,
    sharp corners,boxrule=0pt,width=5.5em,height=1.4em,
    halign=center,valign=center,nobeforeafter,left=0pt,right=0pt,
    top=0pt,bottom=0pt]\small\bfseries Blueprint\end{tcolorbox}}

\newcommand{\verselabel}{%
  \noindent\begin{tcolorbox}[colback=verseheadbg,coltext=white,
    sharp corners,boxrule=0pt,width=3.8em,height=1.4em,
    halign=center,valign=center,nobeforeafter,left=0pt,right=0pt,
    top=0pt,bottom=0pt]\small\bfseries Verse\end{tcolorbox}}

\newcommand{\tcllabel}{%
  \noindent\begin{tcolorbox}[colback=tclheadbg,coltext=white,
    sharp corners,boxrule=0pt,width=2.8em,height=1.4em,
    halign=center,valign=center,nobeforeafter,left=0pt,right=0pt,
    top=0pt,bottom=0pt]\small\bfseries Tcl\end{tcolorbox}}

\newcommand{\jslabel}{%
  \noindent\begin{tcolorbox}[colback=jsheadbg!80!black,coltext=white,
    sharp corners,boxrule=0pt,width=6.0em,height=1.4em,
    halign=center,valign=center,nobeforeafter,left=0pt,right=0pt,
    top=0pt,bottom=0pt]\small\bfseries JavaScript\end{tcolorbox}}

\newcommand{\guilelabel}{%
  \noindent\begin{tcolorbox}[colback=guileheadbg,coltext=white,
    sharp corners,boxrule=0pt,width=3.5em,height=1.4em,
    halign=center,valign=center,nobeforeafter,left=0pt,right=0pt,
    top=0pt,bottom=0pt]\small\bfseries Guile\end{tcolorbox}}

\newcommand{\rhailabel}{%
  \noindent\begin{tcolorbox}[colback=rhaiheadbg,coltext=white,
    sharp corners,boxrule=0pt,width=3.2em,height=1.4em,
    halign=center,valign=center,nobeforeafter,left=0pt,right=0pt,
    top=0pt,bottom=0pt]\small\bfseries Rhai\end{tcolorbox}}

\newcommand{\vbalabel}{%
  \noindent\begin{tcolorbox}[colback=vbaheadbg,coltext=white,
    sharp corners,boxrule=0pt,width=3.0em,height=1.4em,
    halign=center,valign=center,nobeforeafter,left=0pt,right=0pt,
    top=0pt,bottom=0pt]\small\bfseries VBA\end{tcolorbox}}

% ======================================================
\begin{document}
% ======================================================

\begin{titlepage}
\centering
\vspace*{3cm}
{\Huge\bfseries\color{algheadbg} C/C++ 脚本语言集成对比参考手册\par}
\vspace{1cm}
{\LARGE\color{blue!70!black} Lua \textbullet{} Python \textbullet{} AngelScript \textbullet{} ChaiScript \textbullet{} Squirrel\par}
\vspace{0.3cm}
{\LARGE\color{blue!70!black} UE Blueprint \textbullet{} Verse \textbullet{} Tcl \textbullet{} V8/JavaScript \textbullet{} Guile \textbullet{} Rhai \textbullet{} VBA\par}
\vspace{0.3cm}
{\LARGE\color{blue!70!black} C\# 脚本化方案对比\par}
\vfill
{\large\color{gray} \today\par}
\end{titlepage}

\frontmatter

% ─── 前言 ───
\chapter*{前言}
\addcontentsline{toc}{chapter}{前言}

在大型 C/C++ 项目中，频繁的编译--链接--部署循环严重拖慢开发迭代速度。将部分逻辑下沉到脚本语言中，可以在不重新编译宿主程序的前提下实现行为定制、热更新和快速原型验证。

本手册系统梳理了 C/C++ 中主流的脚本语言嵌入方案，涵盖 Lua（及 sol2 绑定库）、Python（CPython API 与 pybind11）、AngelScript、ChaiScript 和 Squirrel。同时收录了游戏工业中的 UE Blueprint 蓝图系统和 Verse 语言，以及工业与学术领域常用的 Tcl、V8/JavaScript、Guile/Scheme 和 Rhai 等嵌入式脚本引擎。每种方案从虚拟机初始化、函数注册与调用、数据交换、类型映射到热重载机制均有详细代码示例。手册单独设立 C\# 部分，介绍 Roslyn Scripting、\Fn{dynamic} 关键字和反射机制如何实现类似的动态执行能力。最后通过综合对比帮助读者在不同场景下做出技术选型。

% ─── 目录 ───
\makeatletter
\IfFileExists{\jobname-manual.toc}{%
  \chapter*{目录}
  \@input{\jobname-manual.toc}%
}{%
  \tableofcontents
}
\makeatother

\mainmatter

% ======================================================
\part{C/C++ 脚本集成方案}
% ======================================================

% ====================================
\chapter{脚本集成概述}\label{ch:overview}
% ====================================

\section{为什么需要脚本集成}\label{sec:why}

C/C++ 以编译型语言著称，编译产出高性能的原生代码，但也意味着每次修改逻辑都需要经历 \textbf{编辑 $\to$ 编译 $\to$ 链接 $\to$ 部署} 的完整流程。对于大型项目（如游戏引擎、嵌入式控制系统、网络服务），一次全量编译可能耗时数分钟甚至更长。脚本集成通过将部分业务逻辑外置到解释执行的脚本语言中，实现：

\begin{itemize}
  \item \textbf{零编译热更新}：修改脚本文件后直接加载执行，无需重新编译宿主程序。
  \item \textbf{快速原型验证}：用脚本语言快速编写和测试算法逻辑，验证后再移植到 C/C++。
  \item \textbf{用户可定制行为}：暴露脚本接口给终端用户或关卡设计师，允许他们编写自定义逻辑。
  \item \textbf{降低门槛}：脚本语言语法更简洁，适合非核心引擎开发人员参与。
\end{itemize}

\section{嵌入式脚本架构模型}\label{sec:arch}

所有嵌入式脚本方案都遵循相似的架构模型：

\begin{enumerate}
  \item \textbf{宿主进程}（Host）：C/C++ 编写的核心程序，负责初始化脚本虚拟机、注册 C/C++ 函数供脚本调用、加载和执行脚本。
  \item \textbf{脚本虚拟机}（VM）：脚本语言的运行时环境，负责解析、编译（或解释）和执行脚本代码。可以是嵌入式 VM（Lua、AngelScript）或独立进程（Python 解释器）。
  \item \textbf{绑定层}（Binding Layer）：连接宿主与 VM 的桥梁，负责类型转换、函数签名映射、异常传播和生命周期管理。
  \item \textbf{脚本文件}（Script）：纯文本文件，包含用脚本语言编写的业务逻辑。
\end{enumerate}

\begin{guideline}[title={选型考量因素}]
\begin{tabularx}{\textwidth}{lX}
\toprule
\textbf{因素} & \textbf{说明} \\
\midrule
运行时体积 & Lua $<$ 250\,KB，Squirrel $\sim$300\,KB，AngelScript $\sim$500\,KB，Python $\sim$30\,MB \\
执行性能 & LuaJIT 接近原生，CPython 最慢，AngelScript 可选 JIT 编译 \\
绑定复杂度 & sol2/pybind11 现代化，CPython API 最底层 \\
生态与库 & Python 最丰富，Lua 有大量游戏开发库 \\
许可证 & Lua MIT，Python PSF，AngelScript zlib，ChaiScript BSD \\
\bottomrule
\end{tabularx}
\end{guideline}

\section{编译开销对比}\label{sec:compile-cost}

脚本集成最直接的收益是减少编译开销。下表展示了典型场景下的对比：

\begin{center}
\begin{tabularx}{\textwidth}{lXX}
\toprule
\textbf{场景} & \textbf{纯 C++ 方案} & \textbf{C++ + 脚本方案} \\
\midrule
修改一行业务逻辑 & 重新编译整个翻译单元 + 链接 & 保存脚本文件，立即生效 \\
新增一个事件处理器 & 编写 C++ 代码 + 编译 + 部署 & 在脚本中新增函数即可 \\
调整 AI 行为参数 & 修改代码或配置文件 + 重启 & 热重载脚本，无需重启 \\
A/B 测试不同策略 & 条件编译或多态 + 重新编译 & 切换脚本文件路径即可 \\
\bottomrule
\end{tabularx}
\end{center}


% ====================================
\chapter{Lua 嵌入}\label{ch:lua}
% ====================================

Lua 是最流行的嵌入式脚本语言，其设计目标就是作为宿主应用程序的扩展语言。Lua 虚拟机极其轻量（默认栈大小仅几十 KB），启动开销微乎其微，且采用 MIT 许可证。

\section{C API 基础}\label{sec:lua-capi}

\subsection{虚拟机生命周期}\label{sec:lua-vm}

Lua 通过 \Fn{lua\_newstate} 或 \Fn{luaL\_newstate} 创建虚拟机实例。每个 VM 拥有独立的栈和全局环境，可在同一进程中运行多个 VM 实现隔离：

\cpplabel
\begin{lstlisting}[style=cppstyle]
#include <lua.hpp>
#include <lauxlib.h>
#include <lualib.h>

int main() {
    // 创建新状态机（使用默认分配器）
    lua_State* L = luaL_newstate();
    if (!L) return 1;

    // 加载标准库（math, string, table, io 等）
    luaL_openlibs(L);

    // 执行脚本文件
    if (luaL_dofile(L, "game_logic.lua") != LUA_OK) {
        const char* err = lua_tostring(L, -1);
        fprintf(stderr, "Lua error: %s\n", err);
        lua_pop(L, 1);
    }

    // 清理
    lua_close(L);
    return 0;
}
\end{lstlisting}

\begin{note}[title={多 VM 隔离}]
每个 \Tp{lua\_State} 完全独立：全局变量、注册表、元表互不干扰。游戏中常见做法是为每个子系统（AI、UI、事件）分配独立 VM，避免脚本间意外耦合。
\end{note}

\subsection{栈操作模型}\label{sec:lua-stack}

Lua C API 的核心概念是 \textbf{虚拟栈}（Virtual Stack）。C 与 Lua 之间的所有数据交换都通过栈进行，栈顶索引为 $-1$，栈底为 $1$：

\cpplabel
\begin{lstlisting}[style=cppstyle]
// 向栈压入值
lua_pushinteger(L, 42);           // 栈: [42]
lua_pushstring(L, "hello");       // 栈: [42, "hello"]
lua_pushboolean(L, 1);            // 栈: [42, "hello", true]

// 读取栈上的值（不移除）
int val = lua_tointeger(L, -3);   // val = 42
const char* s = lua_tostring(L, -2); // s = "hello"

// 弹出 n 个元素
lua_pop(L, 2);                    // 栈: [42]
\end{lstlisting}

\begin{warning}[title={栈平衡}]
每个 C 函数被 Lua 调用时，必须在返回前将栈恢复到调用前的状态（或仅保留返回值）。栈溢出或栈下溢是最常见的 Lua 集成 bug。使用 \Fn{lua\_gettop} 检查栈深度，\Fn{lua\_checkstack} 确保栈有足够空间。
\end{warning}

\section{C 函数注册}\label{sec:lua-func-reg}

\subsection{注册与调用}\label{sec:lua-reg-call}

C 函数必须遵循 \Tp{lua\_CFunction} 签名（\Tp{int (*)(lua\_State*)}），通过栈接收参数、通过栈返回结果：

\cpplabel
\begin{lstlisting}[style=cppstyle]
// C 函数：计算两个向量的点积
static int l_dot_product(lua_State* L) {
    double x1 = luaL_checknumber(L, 1); (*@\textcolor{gray}{// 第1个参数}@*)
    double y1 = luaL_checknumber(L, 2);
    double x2 = luaL_checknumber(L, 3);
    double y2 = luaL_checknumber(L, 4);

    lua_pushnumber(L, x1 * x2 + y1 * y2); (*@\textcolor{gray}{// 返回值}@*)
    return 1; (*@\textcolor{gray}{// 返回1个值}@*)
}

// 注册到全局表
static const luaL_Reg mylib[] = {
    {"dot_product", l_dot_product},
    {NULL, NULL}
};

// 在初始化时注册
luaL_newlib(L, mylib);        (*@\textcolor{gray}{// 创建表并填充函数}@*)
lua_setglobal(L, "mylib");    (*@\textcolor{gray}{// mylib = \{dot\_product = ...\}}@*)
\end{lstlisting}

\lualabel
\begin{lstlisting}[style=luastyle]
-- 在 Lua 脚本中调用注册的 C 函数
local result = mylib.dot_product(1.0, 2.0, 3.0, 4.0)
print("dot product:", result)  -- 输出: 11.0
\end{lstlisting}

\subsection{从 C 调用 Lua 函数}\label{sec:lua-call-lua}

宿主程序也可以调用 Lua 中定义的函数。方法是依次将函数和参数压入栈，然后调用 \Fn{lua\_pcall} 完成调用：

\cpplabel
\begin{lstlisting}[style=cppstyle]
// 调用 Lua 函数: result = on_damage(target, amount)
lua_getglobal(L, "on_damage");     (*@\textcolor{gray}{// 压入函数}@*)
lua_pushstring(L, "player_1");     (*@\textcolor{gray}{// 参数1}@*)
lua_pushinteger(L, 25);            (*@\textcolor{gray}{// 参数2}@*)

if (lua_pcall(L, 2, 1, 0) != LUA_OK) { (*@\textcolor{gray}{// 2参数, 1返回值}@*)
    fprintf(stderr, "error: %s\n",
            lua_tostring(L, -1));
    lua_pop(L, 1);
} else {
    int result = lua_tointeger(L, -1);
    lua_pop(L, 1);
    printf("damage result: %d\n", result);
}
\end{lstlisting}

\section{数据交换}\label{sec:lua-data}

\subsection{表（Table）操作}\label{sec:lua-table}

Lua 的表是唯一的复合数据结构，同时充当数组和字典。C 端通过栈操作读写表：

\cpplabel
\begin{lstlisting}[style=cppstyle]
// 读取 Lua 表: config = {max_hp = 100, speed = 5.5}
lua_getglobal(L, "config");

lua_getfield(L, -1, "max_hp");  (*@\textcolor{gray}{// 从表中取字段}@*)
int max_hp = lua_tointeger(L, -1);
lua_pop(L, 1);

lua_getfield(L, -1, "speed");
double speed = lua_tonumber(L, -1);
lua_pop(L, 1);

lua_pop(L, 1); (*@\textcolor{gray}{// 弹出 config 表}@*)
\end{lstlisting}

\subsection{Userdata：C 对象桥接}\label{sec:lua-userdata}

\Tp{userdata} 允许在 Lua 中持有 C/C++ 对象的指针，通过元表（metatable）为其附加方法：

\cpplabel
\begin{lstlisting}[style=cppstyle]
struct Entity { int hp; float x, y; };

static int l_entity_take_damage(lua_State* L) {
    Entity* e = (Entity*)luaL_checkudata(L, 1, "Entity");
    int dmg = luaL_checkinteger(L, 2);
    e->hp -= dmg;
    lua_pushinteger(L, e->hp);
    return 1;
}

static int l_entity_gc(lua_State* L) {
    Entity* e = (Entity*)luaL_checkudata(L, 1, "Entity");
    e->~Entity(); (*@\textcolor{gray}{// 显式析构}@*)
    return 0;
}

// 注册 Entity 元表
static const luaL_Reg entity_methods[] = {
    {"take_damage", l_entity_take_damage},
    {"__gc",        l_entity_gc},
    {NULL, NULL}
};

void register_entity(lua_State* L) {
    luaL_newmetatable(L, "Entity");
    luaL_setfuncs(L, entity_methods, 0);
    lua_pushvalue(L, -1);
    lua_setfield(L, -2, "__index"); (*@\textcolor{gray}{// mt.\_\_index = mt}@*)
    lua_pop(L, 1);
}
\end{lstlisting}

\lualabel
\begin{lstlisting}[style=luastyle]
-- Lua 侧使用 userdata
local e = create_entity(100, 0.0, 0.0)
e:take_damage(30)
print(e.hp)  -- 70
\end{lstlisting}

\section{LuaJIT 与 FFI}\label{sec:luajit}

LuaJIT 提供了接近原生代码的执行性能，其 FFI（Foreign Function Interface）库允许 Lua 直接调用 C 函数和访问 C 数据结构，无需经过传统的 C API 绑定：

\lualabel
\begin{lstlisting}[style=luastyle]
local ffi = require("ffi")

-- 声明 C 函数签名
ffi.cdef[[
    int printf(const char* fmt, ...);
    typedef struct { int x, y; } Point;
    double distance(Point* a, Point* b);
]]

-- 直接调用 C 函数（无需绑定代码！）
ffi.C.printf("Hello from LuaJIT FFI!\n")

-- 创建 C 结构体实例
local p1 = ffi.new("Point", 0, 0)
local p2 = ffi.new("Point", 3, 4)
local dist = ffi.C.distance(p1, p2)
print("distance:", dist)  -- 5.0
\end{lstlisting}

\begin{bestpractice}[title={LuaJIT FFI 使用建议}]
LuaJIT FFI 的 \Fn{ffi.cdef} 可以直接解析 C 头文件声明，省去手写绑定代码。对于性能敏感的数学计算、数据转换等场景，FFI 调用开销远低于通过 Lua C API 栈传递。但 FFI 不支持 C++ 类、模板和重载——需要先用 \texttt{extern "C"} 包装 C++ 接口。
\end{bestpractice}


% ====================================
\chapter{sol2：现代 C++ Lua 绑定}\label{ch:sol2}
% ====================================

sol2 是一个 header-only 的 C++17 库，将 Lua C API 包装为类型安全的现代 C++ 接口，消除了手写栈操作和类型转换的繁琐代码。

\section{基本用法}\label{sec:sol2-basic}

\sollabel
\begin{lstlisting}[style=cppstyle]
#include <sol/sol.hpp>

int main() {
    sol::state lua;
    lua.open_libraries(sol::lib::base, sol::lib::math,
                       sol::lib::string);

    // 直接设置全局变量
    lua["app_name"] = "MyGame";
    lua["version"]  = 2;

    // 注册 C++ 函数（自动类型转换）
    lua["add"] = [](int a, int b) { return a + b; };

    // 执行脚本
    lua.script(R"(
        print(app_name .. " v" .. version)
        local sum = add(10, 20)
        print("10 + 20 = " .. sum)
    )");

    return 0;
}
\end{lstlisting}

\section{类绑定}\label{sec:sol2-class}

sol2 的类绑定是声明式的，支持构造函数、成员函数、属性和继承：

\sollabel
\begin{lstlisting}[style=cppstyle]
struct Vec2 {
    float x, y;
    Vec2(float x, float y) : x(x), y(y) {}
    float length() const { return std::sqrt(x*x + y*y); }
    Vec2 operator+(const Vec2& o) const {
        return {x + o.x, y + o.y};
    }
};

void register_vec2(sol::state& lua) {
    lua.new_usertype<Vec2>("Vec2",
        sol::constructors<Vec2(float, float)>(),
        "x", &Vec2::x,             (*@\textcolor{gray}{// 可读写属性}@*)
        "y", &Vec2::y,
        "length", &Vec2::length,   (*@\textcolor{gray}{// 只读方法}@*)
        sol::meta_function::addition,
            &Vec2::operator+        (*@\textcolor{gray}{// 运算符重载}@*)
    );
}
\end{lstlisting}

\lualabel
\begin{lstlisting}[style=luastyle]
local v1 = Vec2.new(3, 4)
local v2 = Vec2.new(1, 2)
local v3 = v1 + v2
print(v3.x, v3.y)        -- 4.0  6.0
print(v1:length())        -- 5.0
v1.x = 10                 -- 直接修改属性
\end{lstlisting}

\section{错误处理}\label{sec:sol2-error}

sol2 将 Lua 错误映射为 C++ 异常，支持 \Tp{sol::protected\_function} 进行安全调用：

\sollabel
\begin{lstlisting}[style=cppstyle]
sol::state lua;
lua.open_libraries(sol::lib::base);

// 加载有潜在错误的脚本
auto result = lua.safe_script(R"(
    function risky(x)
        if x < 0 then error("negative!") end
        return math.sqrt(x)
    end
)", sol::script_pass_on_error);

// 安全调用
sol::protected_function risky = lua["risky"];
auto [ok, val] = risky(-1);
if (!ok) {
    sol::error err = val;
    std::cerr << "Lua error: " << err.what() << '\n';
    (*@\textcolor{gray}{// 输出: Lua error: [string "..."]:3: negative!}@*)
}
\end{lstlisting}

\begin{guideline}[title={sol2 vs 原始 C API}]
sol2 将典型绑定代码量减少 60--80\%，代价是编译时间增加（大量模板展开）。对于绑定层较薄的项目（$<$50 个函数），编译时间影响可忽略。对于大型项目，建议将绑定代码隔离在独立的翻译单元中，利用预编译头加速。
\end{guideline}


% ====================================
\chapter{Python 嵌入}\label{ch:python}
% ====================================

Python 是最广泛使用的脚本语言，拥有庞大的标准库和第三方生态。将 Python 嵌入 C/C++ 应用程序可以利用其丰富的功能，但代价是较大的运行时体积（$\sim$30\,MB）和 GIL 带来的多线程限制。

\section{CPython C API}\label{sec:cpython-api}

\subsection{解释器初始化}\label{sec:python-init}

\cpplabel
\begin{lstlisting}[style=cppstyle]
#include <Python.h>

int main() {
    Py_Initialize(); (*@\textcolor{gray}{// 初始化 Python 解释器}@*)

    // 执行 Python 代码
    PyRun_SimpleString(
        "import sys\n"
        "print(f'Python {sys.version}')"
    );

    // 导入并调用模块函数
    PyObject* module = PyImport_ImportModule("math");
    if (module) {
        PyObject* func = PyObject_GetAttrString(
            module, "sqrt");
        PyObject* args = PyTuple_Pack(
            1, PyFloat_FromDouble(144.0));
        PyObject* result = PyObject_Call(func, args, NULL);
        if (result) {
            printf("sqrt(144) = %f\n",
                   PyFloat_AsDouble(result));
            Py_DECREF(result);
        }
        Py_DECREF(args);
        Py_DECREF(func);
        Py_DECREF(module);
    }

    Py_Finalize();
    return 0;
}
\end{lstlisting}

\begin{warning}[title={引用计数}]
CPython API 使用手动引用计数管理对象生命周期。\Fn{Py\_DECREF} 遗漏会导致内存泄漏，多余的 \Fn{Py\_DECREF} 会导致崩溃。这是 CPython API 最大的痛点——也是 pybind11 存在的理由。
\end{warning}

\subsection{GIL 与多线程}\label{sec:python-gil}

Python 的全局解释器锁（GIL）确保同一时刻只有一个线程执行 Python 字节码。在多线程 C++ 应用中嵌入 Python 时：

\cpplabel
\begin{lstlisting}[style=cppstyle]
// 主线程初始化后释放 GIL
Py_Initialize();
PyEval_SaveThread(); (*@\textcolor{gray}{// 释放 GIL，允许其他线程获取}@*)

// 工作线程需要 Python 时重新获取 GIL
void worker_thread() {
    PyGILState_STATE gstate = PyGILState_Ensure();

    PyRun_SimpleString("print('from thread')");

    PyGILState_Release(gstate); (*@\textcolor{gray}{// 用完释放}@*)
}

// 程序退出前
PyGILState_Ensure();
Py_Finalize();
\end{lstlisting}

\section{pybind11}\label{sec:pybind11}

pybind11 是 Python 绑定的现代 C++ 库（header-only，C++11），其设计理念与 sol2 之于 Lua 类似——用模板和 RAII 消除手动引用计数和栈操作。

\pybindlabel
\begin{lstlisting}[style=cppstyle]
#include <pybind11/pybind11.h>
#include <pybind11/stl.h>
namespace py = pybind11;

struct Config {
    std::string name;
    int max_retries = 3;
    std::vector<std::string> plugins;
};

// 绑定模块定义
PYBIND11_EMBEDDED_MODULE(myapp, m) {
    py::class_<Config>(m, "Config")
        .def(py::init<>())
        .def_readwrite("name", &Config::name)
        .def_readwrite("max_retries",
                        &Config::max_retries)
        .def_readwrite("plugins", &Config::plugins)
        .def("__repr__", [](const Config& c) {
            return "<Config '" + c.name + "'>";
        });

    m.def("load_config", [](const std::string& path) {
        Config cfg;
        cfg.name = path; (*@\textcolor{gray}{// 简化示例}@*)
        return cfg;
    });
}
\end{lstlisting}

\cpplabel
\begin{lstlisting}[style=cppstyle]
#include <pybind11/embed.h>
namespace py = pybind11;

int main() {
    py::scoped_interpreter guard{}; (*@\textcolor{gray}{// RAII 管理解释器生命周期}@*)

    py::exec(R"(
        import myapp
        cfg = myapp.load_config("app.yaml")
        cfg.max_retries = 5
        cfg.plugins = ["auth", "logging"]
        print(cfg)
    )");

    // C++ 侧也可以访问 Python 对象
    auto cfg = py::module_::import("myapp")
                   .attr("load_config")("test.yaml");
    std::string name = cfg.attr("name").cast<std::string>();

    return 0;
} (*@\textcolor{gray}{// guard 析构时自动 Py\_Finalize()}@*)
\end{lstlisting}

\begin{bestpractice}[title={pybind11 与 sol2 的相似性}]
pybind11 和 sol2 采用了几乎相同的设计理念：利用 C++ 模板推导和 RAII 将底层 C API 的繁琐操作（栈管理、引用计数、类型检查）封装为类型安全的声明式接口。掌握其中一个后，学习另一个的成本很低。
\end{bestpractice}


% ====================================
\chapter{AngelScript}\label{ch:angelscript}
% ====================================

AngelScript 是专门为游戏和实时应用设计的嵌入式脚本语言。它的语法与 C/C++ 非常相似（静态类型、支持类和接口），学习曲线对 C++ 开发者极为平缓。采用 zlib 许可证。

\section{引擎与模块}\label{sec:as-engine}

\angelscriptlabel
\begin{lstlisting}[style=cppstyle]
#include <angelscript.h>
#include <cassert>

void message_callback(const asSMessageInfo* msg,
                      void* param) {
    printf("%s (%d,%d): %s\n",
           msg->section, msg->row, msg->col,
           msg->message);
}

int main() {
    asIScriptEngine* engine = asCreateScriptEngine();

    // 设置消息回调（编译错误/警告）
    engine->SetMessageCallback(
        asFUNCTION(message_callback), nullptr,
        asCALL_CDECL);

    // 注册 C++ 函数供脚本调用
    engine->RegisterGlobalFunction(
        "void print(const string &in msg)",
        asFUNCTION([](const std::string& s) {
            printf("%s\n", s.c_str());
        }),
        asCALL_CDECL);

    // 编译并执行脚本
    asIScriptModule* mod =
        engine->GetModule("game", asGM_ALWAYS_CREATE);
    mod->AddScriptSection("main",
        R"(void OnStart() {
            print("AngelScript running!");
        })");
    mod->Build();

    // 调用脚本函数
    asIScriptContext* ctx = engine->CreateContext();
    asIScriptFunction* func =
        mod->GetFunctionByDecl("void OnStart()");
    ctx->Prepare(func);
    ctx->Execute();
    ctx->Release();

    engine->ShutDownAndRelease();
    return 0;
}
\end{lstlisting}

\section{类型注册}\label{sec:as-type-reg}

AngelScript 的类型注册是显式的，需要逐一注册属性和方法：

\cpplabel
\begin{lstlisting}[style=cppstyle]
struct Vec3 {
    float x, y, z;
    Vec3() : x(0), y(0), z(0) {}
    Vec3(float x, float y, float z)
        : x(x), y(y), z(z) {}
    float Length() const {
        return sqrtf(x*x + y*y + z*z);
    }
};

void RegisterVec3(asIScriptEngine* engine) {
    // 注册值类型
    engine->RegisterObjectType("Vec3",
        sizeof(Vec3), asOBJ_VALUE | asOBJ_POD);

    // 构造函数
    engine->RegisterObjectBehaviour("Vec3",
        asBEHAVE_CONSTRUCT,
        "void f(float x, float y, float z)",
        asFUNCTIONPR(
            +[](void* mem, float x, float y, float z) {
                new(mem) Vec3(x, y, z);
            },
            void(void*, float, float, float)),
        asCALL_CDECL_OBJLAST);

    // 成员属性
    engine->RegisterObjectProperty("Vec3",
        "float x", asOFFSET(Vec3, x));
    engine->RegisterObjectProperty("Vec3",
        "float y", asOFFSET(Vec3, y));
    engine->RegisterObjectProperty("Vec3",
        "float z", asOFFSET(Vec3, z));

    // 成员方法
    engine->RegisterObjectMethod("Vec3",
        "float Length() const",
        asMETHOD(Vec3, Length), asCALL_THISCALL);
}
\end{lstlisting}

\begin{compare}[title={AngelScript vs Lua}]
AngelScript 的优势在于：(1) 静态类型，编译时即可发现类型错误；(2) 语法接近 C++，开发者无需学习新范式；(3) 原生支持类和接口继承。劣势在于：(1) 类型注册比 sol2 更繁琐；(2) 社区和生态远小于 Lua/Python；(3) 没有 JIT 编译（但有字节码预编译）。
\end{compare}

\section{字节码序列化}\label{sec:as-bytecode}

AngelScript 支持将编译后的脚本保存为字节码，运行时直接加载字节码跳过编译步骤：

\cpplabel
\begin{lstlisting}[style=cppstyle]
// 保存字节码
asIScriptModule* mod = engine->GetModule("game");
class CBytecodeStream : public asIByteCodeStream {
    std::vector<char> data; (*@\textcolor{gray}{// 简化实现}@*)
    // ... 实现 Write, Read, GetSize 接口 ...
};
CBytecodeStream stream;
mod->SaveByteCode(&stream);
// 将 stream.data 写入磁盘文件

// 加载字节码（在另一台机器或下次启动时）
asIScriptModule* mod2 =
    engine->GetModule("game", asGM_ALWAYS_CREATE);
mod2->LoadByteCode(&stream);
\end{lstlisting}


% ====================================
\chapter{其他脚本引擎}\label{ch:others}
% ====================================

\section{ChaiScript}\label{sec:chaiscript}

ChaiScript 是一个 header-only 的 C++ 嵌入式脚本语言，其最大特色是 \textbf{与 C++ 语法几乎一致}——同一个源文件中可以混合使用 C++ 和 ChaiScript 代码。

\chaiscriptlabel
\begin{lstlisting}[style=cppstyle]
#include <chaiscript/chaiscript.hpp>

struct Player {
    std::string name;
    int hp = 100;
    void take_damage(int dmg) { hp -= dmg; }
};

int main() {
    chaiscript::ChaiScript chai;

    // 注册类型（语法接近 C++ 声明）
    chai.add(chaiscript::user_type<Player>(), "Player");
    chai.add(chaiscript::constructor<Player()>(),
             "Player");
    chai.add(chaiscript::fun(&Player::name), "name");
    chai.add(chaiscript::fun(&Player::hp), "hp");
    chai.add(chaiscript::fun(&Player::take_damage),
             "take_damage");

    // 注册自由函数
    chai.add(chaiscript::fun(
        [](Player& p, int bonus) {
            p.hp += bonus;
        }), "heal");

    // 执行脚本（语法几乎就是 C++）
    chai.eval(R"(
        var p = Player()
        p.name = "Hero"
        p.take_damage(30)
        heal(p, 10)
        print("HP: " + to_string(p.hp))  (*@\textcolor{gray}{// 80}@*)
    )");

    return 0;
}
\end{lstlisting}

\begin{note}[title={ChaiScript 的编译代价}]
ChaiScript 大量使用 C++ 模板和 \Tp{std::function} 包装，导致包含其头文件的翻译单元编译时间显著增加。建议将 ChaiScript 相关代码隔离到独立的 .cpp 文件中，避免污染其他模块的头文件依赖。
\end{note}

\section{Squirrel}\label{sec:squirrel}

Squirrel 是专为游戏开发设计的嵌入式脚本语言，语法借鉴了 C、Lua 和 JavaScript。Valve 的 Source 2 引擎曾使用 Squirrel 作为脚本语言。

\squirrellabel
\begin{lstlisting}[style=cppstyle]
#include <squirrel.h>
#include <sqstdio.h>
#include <sqstdaux.h>

void print_func(HSQUIRRELVM v, const SQChar* s, ...) {
    va_list args;
    va_start(args, s);
    vprintf(s, args);
    va_end(args);
}

int main() {
    HSQUIRRELVM v = sq_open(1024); (*@\textcolor{gray}{// 栈大小}@*)
    sqstd_seterrorhandlers(v);
    sq_setprintfunc(v, print_func, print_func);

    // 编译并执行脚本
    sq_compile(v, /*readfn=*/nullptr,
        R"(
            class Animal {
                name = null;
                constructor(n) { name = n; }
                function speak() {
                    return name + " makes a sound";
                }
            }
            local dog = Animal("Rex");
            print(dog.speak());
        )",
        "script", _SC("script"), SQTrue);
    sq_push(v, -1); (*@\textcolor{gray}{// 压入编译结果（闭包）}@*)
    sq_pushroottable(v); (*@\textcolor{gray}{// this = root table}@*)
    sq_call(v, 1, SQFalse, SQTrue);
    sq_pop(v, 1);

    sq_close(v);
    return 0;
}
\end{lstlisting}

\section{方案选型矩阵}\label{sec:matrix}

\begin{compare}[title={传统文本脚本引擎综合对比}]
\begin{tabularx}{\textwidth}{lXXXXX}
\toprule
\textbf{特性} & \textbf{Lua} & \textbf{Python} & \textbf{Angel\seqsplit{Script}} & \textbf{Chai\seqsplit{Script}} & \textbf{Squirrel} \\
\midrule
体积 & $\sim$250\,KB & $\sim$30\,MB & $\sim$500\,KB & header-only & $\sim$300\,KB \\
类型系统 & 动态 & 动态 & 静态 & 动态 & 动态 \\
语法风格 & 极简 & Pythonic & C-like & C++-like & C/JS-like \\
JIT & LuaJIT & 无 & 无 & 无 & 无 \\
绑定库 & sol2 & pybind11 & 原生API & header-only & 原生API \\
多线程 & 多VM & GIL限制 & 多引擎 & 单线程 & 多VM \\
许可证 & MIT & PSF & zlib & BSD & MIT \\
\bottomrule
\end{tabularx}
\end{compare}

\begin{compare}[title={游戏引擎脚本与工业/办公脚本对比}]
{\small
\begin{tabularx}{\textwidth}{lXXXXX}
\toprule
\textbf{特性} & \textbf{Blueprint} & \textbf{Verse} & \textbf{Tcl} & \textbf{Guile} & \textbf{VBA} \\
\midrule
形态 & 可视化/\seqsplit{Verse}前端 & 文本（函数式） & 文本（命令式） & 文本（Lisp） & 文本（BASIC） \\
宿主 & UE 引擎 & UEFN / UE6+ & 任意 C/C++ & GNU 项目 & Office / CAD \\
类型系统 & 动态 & 静态+逻辑 & 动态 & 动态 & COM VARIANT \\
热更新 & 保存即生效 & 保存即生效 & 重新 Eval & 重新 Eval & 即时执行 \\
典型场景 & 关卡/UI/AI & 创作/脚本 & EDA/CAD & 学术/DSL & 报表/参数化 \\
学习门槛 & 低 & 中 & 低 & 高（Lisp） & 低（BASIC） \\
\bottomrule
\end{tabularx}}
\end{compare}


% ====================================
\chapter{UE Blueprint 蓝图系统}\label{ch:blueprint}
% ====================================

Unreal Engine 的 Blueprint（蓝图）是游戏行业中最成功的可视化脚本系统。它允许设计师和开发者通过拖拽节点和连线来构建游戏逻辑，无需编写传统 C++ 代码即可实现复杂行为，从根本上消除了编译等待。不过，Epic Games 已在 UE6 路线图中明确表示将逐步用 Verse 取代蓝图作为主要的用户侧脚本语言——这意味着蓝图将过渡为 Verse 的可视化前端，而非独立的脚本系统。

\section{蓝图架构}\label{sec:bp-arch}

\blueprintlabel
\begin{guideline}[title={蓝图核心概念}]
\begin{itemize}
  \item \textbf{节点图}（Node Graph）：蓝图由节点（Node）和连线（Wire）组成的有向图。每个节点对应一个函数调用或流程控制操作，连线定义执行顺序和数据流。
  \item \textbf{事件图}（Event Graph）：蓝图的主逻辑区域，以事件节点（如 \Tp{BeginPlay}、\Tp{Tick}）为入口驱动执行。
  \item \textbf{蓝图类型}：Actor Blueprint（场景对象）、Widget Blueprint（UI）、Animation Blueprint（动画状态机）、Behavior Tree（AI 行为树）等。
  \item \textbf{编译产物}：蓝图在编辑器中保存时自动编译为字节码（Blueprint Bytecode），运行时由 UE 虚拟机解释执行。发布时可选择 Nativize（编译为 C++）以获得接近原生的性能。
\end{itemize}
\end{guideline}

\section{蓝图与 C++ 的协作模式}\label{sec:bp-cpp}

在实际项目中，蓝图与 C++ 形成互补的分层架构：

\begin{center}
\begin{tabularx}{\textwidth}{lX}
\toprule
\textbf{层级} & \textbf{职责} \\
\midrule
C++ 基类 & 核心数据结构、高性能算法、网络同步、内存管理等引擎级逻辑 \\
蓝图子类 & 继承 C++ 基类，用可视化节点组合关卡逻辑、角色行为、UI 交互等 \\
蓝图接口 & 定义跨蓝图的通信协议（\Tp{BPI\_XXX}），实现松耦合 \\
\bottomrule
\end{tabularx}
\end{center}

\cpplabel
\begin{lstlisting}[style=cppstyle]
// C++ 基类：暴露属性和函数给蓝图
UCLASS(Blueprintable)
class AMyCharacter : public ACharacter {
    GENERATED_BODY()

    UPROPERTY(EditAnywhere, BlueprintReadWrite,
              Category = "Stats")
    float MaxHealth = 100.0f;

    UPROPERTY(VisibleAnywhere, BlueprintReadOnly,
              Category = "Stats")
    float CurrentHealth;

    // 可在蓝图中调用的 C++ 函数
    UFUNCTION(BlueprintCallable, Category = "Combat")
    void TakeDamage(float Amount);

    // 蓝图可覆写的 C++ 事件
    UFUNCTION(BlueprintImplementableEvent)
    void OnDeath();

protected:
    virtual void BeginPlay() override;
};
\end{lstlisting}

\cpplabel
\begin{lstlisting}[style=cppstyle]
void AMyCharacter::TakeDamage(float Amount) {
    CurrentHealth -= Amount;
    if (CurrentHealth <= 0.0f) {
        OnDeath(); (*@\textcolor{gray}{// 调用蓝图中实现的事件}@*)
    }
}
\end{lstlisting}

蓝图侧通过继承 \Tp{AMyCharacter}，在事件图中用可视化节点实现 \Tp{OnDeath} 逻辑（播放死亡动画、掉落物品、通知 AI 等），无需编译 C++。

\section{蓝图函数库与纯蓝图类}\label{sec:bp-lib}

\textbf{蓝图函数库}（Blueprint Function Library）提供全局可用的纯静态函数，设计师可以在任意蓝图中调用：

\cpplabel
\begin{lstlisting}[style=cppstyle]
UCLASS()
class UGameMathLib : public UBlueprintFunctionLibrary {
    GENERATED_BODY()

    // 纯函数（无副作用，蓝图自动优化）
    UFUNCTION(BlueprintPure, Category = "Math")
    static float CalculateDamage(float BaseDamage,
        float Armor, float Resistance) {
        return BaseDamage
             * (100.0f / (100.0f + Armor))
             * (1.0f - Resistance);
    }

    // 带输出参数的函数
    UFUNCTION(BlueprintCallable, Category = "Util")
    static bool FindNearestEnemy(const FVector& Pos,
        float Radius, AActor*& OutEnemy);
};
\end{lstlisting}

\section{蓝图的编译与性能}\label{sec:bp-perf}

\begin{compare}[title={蓝图 vs C++ 性能对比}]
\begin{tabularx}{\textwidth}{lXX}
\toprule
\textbf{指标} & \textbf{蓝图} & \textbf{C++} \\
\midrule
编译时间 & 保存即编译（$\sim$秒级） & 分钟级（大型项目） \\
运行时性能 & 字节码解释执行，约慢 5--10$\times$ & 原生机器码 \\
热更新 & 编辑器内实时预览 & 需要重新编译部署 \\
调试 & 可视化断点、节点高亮 & 传统调试器 \\
适合场景 & UI、关卡逻辑、快速原型 & 核心系统、高频计算 \\
\bottomrule
\end{tabularx}
\end{compare}

\begin{bestpractice}[title={蓝图 $\to$ Verse/C++ 迁移路径（UE6+）}]
项目初期用蓝图快速验证玩法，性能瓶颈确定后将热点蓝图 \textbf{Nativize} 或手动迁移为 C++。UE 的 \Fn{Blueprint Nativize} 工具可将蓝图自动转换为等价的 C++ 代码并编译，性能提升可达 10$\times$ 以上。

在 UE6+ 路线图中，Epic 计划将 Verse 定位为用户侧脚本的继承者，蓝图将逐步退居为 Verse 的可视化编辑器前端。长期来看，新项目应优先采用 Verse 编写脚本逻辑，蓝图保留用于可视化编辑和快速原型；C++ 继续承担核心系统（物理、AI、网络、渲染）。
\end{bestpractice}


% ====================================
\chapter{Verse 语言}\label{ch:verse}
% ====================================

Verse 是 Epic Games 为 UEFN（Unreal Editor for Fortnite）和未来元宇宙平台设计的新编程语言。它结合了函数式编程、逻辑编程和面向对象编程的特性，目标是成为"元宇宙的 JavaScript"。

\section{语言特性}\label{sec:verse-features}

\verselabel
\begin{guideline}[title={Verse 核心设计}]
\begin{itemize}
  \item \textbf{强类型 + 类型推导}：所有表达式都有静态类型，但大量使用类型推导减少注解。
  \item \textbf{逻辑编程}：支持失败模式（Failure Pattern），表达式可以"失败"（fail）而非抛出异常，类似 Prolog 的回溯机制。
  \item \textbf{时间性}（Temporal）：内建 \Tp{suspend} 和 \Tp{race} 原语，天然支持异步操作和并发。
  \item \textbf{确定性}：在 UEFN 中，Verse 脚本在确定性模拟（Deterministic Simulation）中运行，相同输入保证相同输出。
\end{itemize}
\end{guideline}

\section{Verse 代码示例}\label{sec:verse-examples}

\verselabel
\begin{lstlisting}[style=cppstyle]
# 设备类（UEFN 中的核心概念）
hello_world_device := class(creative_device):

    # OnBegin 在设备激活时自动调用
    OnBegin<override>()<suspends>:void =
        Print("Hello from Verse!")
        # 并发执行多个任务
        spawn:
            CountdownTimer(10)
        spawn:
            MonitorPlayers()

    # 使用 Verse 的失败模式处理可能失败的操作
    CountdownTimer(Seconds:int)<suspends>:void =
        for (I := 0..Seconds):
            Sleep(1.0)
            if (Player := GetPlayspace()
                .GetPlayers()[0]):  (*@\textcolor{gray}{// 失败则跳过}@*)
                Player.SendMessage(
                    "{Seconds - I} seconds left!")

    # 函数式风格：map/filter 链式操作
    GetAlivePlayers():[]player =
        AllPlayers := GetPlayspace().GetPlayers()
        for:
            P := AllPlayers
            IsAlive(P)?    (*@\textcolor{gray}{// 过滤条件}@*)
        yield: P
\end{lstlisting}

\section{Verse 与 UE C++ 的关系}\label{sec:verse-cpp}

\begin{compare}[title={Verse vs Blueprint vs C++ 在 UE 生态中的定位（含 UE6+ 规划）}]
\begin{tabularx}{\textwidth}{lXXX}
\toprule
\textbf{维度} & \textbf{C++} & \textbf{Blueprint} & \textbf{Verse} \\
\midrule
目标用户 & 引擎开发者 & 设计师（UE5） & 脚本开发者 \\
运行环境 & 原生 UE 引擎 & UE 引擎（UE5） & UEFN (UE5) / UE 引擎 (UE6+) \\
UE6+ 定位 & 核心系统 & 退居 Verse 前端 & 用户侧脚本主力 \\
类型系统 & 静态（C++） & 动态（字节码） & 静态 + 逻辑编程 \\
异步模型 & 手动协程/多线程 & 延迟节点 & 内建 suspend/race \\
热更新 & 需要重新编译 & 保存即生效 & 保存即生效 \\
性能 & 最高 & 中等 & 中等（VM 解释执行） \\
\bottomrule
\end{tabularx}
\end{compare}

\begin{note}[title={Verse 的当前状态与 UE6+ 规划}]
Verse 目前仅在 UEFN 环境中可用，无法在标准 Unreal Engine 项目中使用。但 Epic Games 在 UE6 路线图中已将 Verse 定位为蓝图的继承者：未来 Verse 将直接嵌入 UE 引擎运行时，支持在标准 UE 项目中使用，并逐步覆盖蓝图当前的所有用例（关卡逻辑、UI、AI 行为树等）。Verse 目前不提供 FFI 或与 C++ 的直接互操作——与引擎的交互通过 UEFN Creative API 进行，UE6 预计将引入原生 C++ 互操作层。
\end{note}


% ====================================
\chapter{工业与学术脚本嵌入}\label{ch:industrial}
% ====================================

除游戏领域外，大量工业工具（EDA、CAD、科学计算）和学术软件也采用嵌入式脚本实现可扩展性。本章介绍在这些领域中最具代表性的脚本引擎。

\section{Tcl：EDA 与工业自动化的标准脚本}\label{sec:tcl}

Tcl（Tool Command Language）诞生于 1988 年，是 EDA 工具的事实标准脚本语言。Cadence、Synopsys、Mentor 等芯片设计工具链均以 Tcl 作为主要用户接口。

\tcllabel
\begin{lstlisting}[style=cppstyle]
#include <tcl.h>
#include <cstdio>

// 注册给 Tcl 的 C 函数
static int SquareCmd(ClientData, Tcl_Interp* interp,
                     int objc, Tcl_Obj* const objv[]) {
    if (objc != 2) {
        Tcl_WrongNumArgs(interp, 1, objv, "number");
        return TCL_ERROR;
    }
    double val;
    if (Tcl_GetDoubleFromObj(interp, objv[1], &val)
            != TCL_OK) {
        return TCL_ERROR;
    }
    Tcl_SetObjResult(interp,
        Tcl_NewDoubleObj(val * val));
    return TCL_OK;
}

int main() {
    Tcl_Interp* interp = Tcl_CreateInterp();
    Tcl_Init(interp);

    // 注册 C 命令
    Tcl_CreateObjCommand(interp, "square",
        SquareCmd, nullptr, nullptr);

    // 执行 Tcl 脚本
    Tcl_Eval(interp, R"(
        set result [square 7.5]
        puts "7.5 squared = $result"
        # 输出: 7.5 squared = 56.25
    )");

    Tcl_DeleteInterp(interp);
    return 0;
}
\end{lstlisting}

\begin{note}[title={Tcl 的工业地位}]
Tcl 在 EDA 领域的统治地位源于历史惯性：早期的 Cadence SKILL 语言和 Synopsys Design Compiler 都采用 Tcl 作为脚本接口，积累了海量的 Tcl 脚本资产。新项目选择脚本语言时，如果需要与现有 EDA 工具链集成，Tcl 通常是唯一选择。
\end{note}

\section{V8/JavaScript：嵌入浏览器引擎}\label{sec:v8}

V8 是 Google 开发的高性能 JavaScript 引擎，被 Chrome、Node.js、Electron 和大量桌面应用采用。将 V8 嵌入 C++ 应用可以让用户用 JavaScript 编写扩展逻辑。

\jslabel
\begin{lstlisting}[style=cppstyle]
#include <v8.h>
#include <libplatform/libplatform.h>

using namespace v8;

int main() {
    // 初始化 V8
    V8::InitializeICUDefaultLocation("v8_data");
    auto platform = platform::NewDefaultPlatform();
    V8::InitializePlatform(platform.get());
    V8::Initialize();

    Isolate::CreateParams params;
    params.array_buffer_allocator =
        ArrayBuffer::Allocator::NewDefaultAllocator();
    Isolate* isolate = Isolate::New(params);

    {
        Isolate::Scope isolate_scope(isolate);
        HandleScope handle_scope(isolate);
        Local<Context> context =
            Context::New(isolate);
        Context::Scope context_scope(context);

        // 注册 C++ 函数到全局对象
        Local<Object> global = context->Global();
        global->Set(context,
            String::NewFromUtf8Literal(isolate,
                "nativeAdd"),
            Function::New(context,
                [](const FunctionCallbackInfo<
                        Value>& args) {
                    double a = args[0]->NumberValue(
                        args.GetIsolate()->
                        GetCurrentContext()).FromJust();
                    double b = args[1]->NumberValue(
                        args.GetIsolate()->
                        GetCurrentContext()).FromJust();
                    args.GetReturnValue().Set(a + b);
                }).ToLocalChecked()
        ).FromJust();

        // 执行 JavaScript
        const char* src = R"(
            const sum = nativeAdd(100, 200);
            `Result: ${sum}`
        )";
        Local<Script> script =
            Script::Compile(context,
                String::NewFromUtf8Literal(
                    isolate, src))
            .ToLocalChecked();
        Local<Value> result =
            script->Run(context).ToLocalChecked();
        // result = "Result: 300"
    }

    isolate->Dispose();
    V8::Dispose();
    V8::DisposePlatform();
    return 0;
}
\end{lstlisting}

\begin{warning}[title={V8 的嵌入代价}]
V8 的运行时体积约 20--30\,MB，且 API 极其庞大复杂（版本迭代频繁，API 不稳定）。对于只需要轻量脚本能力的项目，Duktape（$\sim$250\,KB，稳定 API）或 QuickJS（$\sim$500\,KB，ES2023 支持）是更实际的嵌入式 JavaScript 引擎选择。
\end{warning}

\section{Guile/Scheme：GNU 扩展语言}\label{sec:guile}

GNU Guile 是 Scheme（Lisp 方言）的一种实现，被设计为 GNU 项目的官方扩展语言。在学术和自由软件社区中，Guile 被广泛用于 Emacs、GnuCash、GNU Artanis 等项目。

\guilelabel
\begin{lstlisting}[style=cppstyle]
#include <libguile.h>

// 注册给 Guile 的 C 函数
static SCM c_multiply(SCM a, SCM b) {
    return scm_product(a, b);
}

static void* register_functions(void* data) {
    scm_c_define_gsubr("c-multiply", 2, 0, 0,
                        (scm_t_subr)c_multiply);
    return nullptr;
}

static void* run_script(void* data) {
    scm_c_eval_string(
        "(define result (c-multiply 6 7))"
        "(format #t \"6 * 7 = ~a~%\" result)"
    ); (*@\textcolor{gray}{// 输出: 6 * 7 = 42}@*)
    return nullptr;
}

int main() {
    scm_with_guile(register_functions, nullptr);
    scm_with_guile(run_script, nullptr);
    return 0;
}
\end{lstlisting}

\begin{compare}[title={Guile 的学术价值}]
Guile/Scheme 的核心价值在于宏系统（\Tp{syntax-rules}、\Tp{syntax-case}）——允许用户在语言层面定义新的语法结构，这是其他嵌入式脚本语言都做不到的。在编译器研究、DSL（领域特定语言）设计和函数式编程教学中，Guile 的宏能力具有不可替代的价值。但 Lisp 的括号语法对大多数终端用户来说学习曲线陡峭。
\end{compare}

\section{Rhai：Rust 生态的嵌入式脚本}\label{sec:rhai}

Rhai 是专为 Rust 设计的嵌入式脚本语言，语法融合了 Rust、JavaScript 与 \seqsplit{RustScript} 的风格。
虽然原生服务于 Rust 生态，也可通过 FFI 桥接嵌入 C++ 项目。

\rhailabel
\begin{lstlisting}[style=cppstyle]
// Rhai 脚本语法示例（非 C++ 代码）
// ─── game_logic.rhai ───
//
// struct Player { name, hp, level }
//
// fn take_damage(self, amount) {
//     self.hp -= amount;
//     if self.hp <= 0 {
//         self.hp = 0;
//         throw "Player died!";
//     }
// }
//
// let player = Player {
//     name: "Hero",
//     hp: 100,
//     level: 1,
// };
//
// player.take_damage(30);
// print(`HP: ${player.hp}`);  // HP: 70
\end{lstlisting}

\begin{compare}[title={Rhai 的定位}]
Rhai 的 API 设计深受 sol2 启发，类型注册使用声明式宏。如果你的项目使用 Rust（而非 C++），Rhai 是目前最成熟的嵌入式脚本选择。对于 C++ 项目，Rhai 需要通过 Rust FFI（\Fn{extern "C"} 包装）桥接，增加了集成复杂度，通常不如直接使用 Lua/sol2 便捷。
\end{compare}

\section{VBA：Office 与 CAD 的自动化标准}\label{sec:vba}

Visual Basic for Applications（VBA）是 Microsoft 推出的嵌入式脚本语言，内置于 Office 套件（Word、Excel、Access、PowerPoint）和大量第三方 CAD/CAE 软件（AutoCAD、SolidWorks、CATIA）中。VBA 是办公自动化和工程设计领域事实上的标准脚本方案。

\vbalabel
\begin{guideline}[title={VBA 嵌入架构}]
\begin{itemize}
  \item \textbf{宿主模型}：VBA 运行在宿主应用程序内嵌的 VBE（Visual Basic Editor）中。宿主通过 COM（Component Object Model）向 VBA 暴露对象模型（Object Model），VBA 代码通过操作这些 COM 对象来控制宿主行为。
  \item \textbf{对象模型}：每个宿主定义自己的对象层次结构。Excel 的 \Tp{Application $\to$ Workbook $\to$ Worksheet $\to$ Range}；AutoCAD 的 \Tp{AcadApplication $\to$ AcadDocument $\to$ ModelSpace}。
  \item \textbf{分发机制}：VBA 使用后期绑定（\Tp{IDispatch}）调用 COM 对象方法，允许脚本在运行时动态发现和调用宿主提供的接口——无需编译时类型信息。
\end{itemize}
\end{guideline}

\subsection{Excel VBA 自动化}\label{sec:vba-excel}

\begin{lstlisting}[style=csharpstyle]
' Excel VBA: 批量生成报表
Sub GenerateReport()
    Dim ws As Worksheet
    Set ws = ThisWorkbook.Sheets("Data")

    Dim lastRow As Long
    lastRow = ws.Cells(ws.Rows.Count, 1).End(xlUp).Row

    ' 遍历数据行，计算汇总
    Dim total As Double: total = 0
    Dim i As Long
    For i = 2 To lastRow
        total = total + ws.Cells(i, 3).Value
    Next i

    ' 写入汇总行
    ws.Cells(lastRow + 2, 2).Value = "Total:"
    ws.Cells(lastRow + 2, 3).Value = total

    ' 格式化和图表生成
    Dim chartObj As ChartObject
    Set chartObj = ws.ChartObjects.Add( _
        Left:=400, Top:=10, _
        Width:=300, Height:=200)
    chartObj.Chart.SetSourceData _
        Source:=ws.Range("A1:C" & lastRow)
    chartObj.Chart.ChartType = xlColumnClustered
End Sub
\end{lstlisting}

\subsection{AutoCAD VBA 自动化}\label{sec:vba-autocad}

\begin{lstlisting}[style=csharpstyle]
' AutoCAD VBA: 批量绘制零件图
Sub DrawBoltPattern()
    Dim doc As AcadDocument
    Set doc = ThisDrawing

    Dim ms As AcadModelSpace
    Set ms = doc.ModelSpace

    ' 参数化设计：螺栓孔圆周阵列
    Dim centerX As Double: centerX = 0#
    Dim centerY As Double: centerY = 0#
    Dim radius As Double: radius = 50#
    Dim boltCount As Integer: boltCount = 8
    Dim boltRadius As Double: boltRadius = 5#

    Dim angle As Double
    Dim i As Integer
    For i = 0 To boltCount - 1
        angle = 2 * 3.14159265 * i / boltCount
        Dim pt(0 To 2) As Double
        pt(0) = centerX + radius * Cos(angle)
        pt(1) = centerY + radius * Sin(angle)
        pt(2) = 0#
        ms.AddCircle pt, boltRadius
    Next i

    ' 添加基圆（参考线）
    ms.AddCircle Array(centerX, centerY, 0#), radius

    doc.Regen acActiveViewport
End Sub
\end{lstlisting}

\subsection{SolidWorks VBA 宏}\label{sec:vba-solidworks}

\begin{lstlisting}[style=csharpstyle]
' SolidWorks VBA: 参数化零件建模
Dim swApp As SldWorks.SldWorks
Dim swModel As SldWorks.ModelDoc2

Sub CreateParametricBox()
    Set swApp = Application.SldWorks
    Set swModel = swApp.ActiveDoc

    ' 读取自定义属性中的尺寸参数
    Dim width As Double, height As Double, depth As Double
    width = swModel.GetCustomInfoValue("", "Width")
    height = swModel.GetCustomInfoValue("", "Height")
    depth = swModel.GetCustomInfoValue("", "Depth")

    ' 在前视基准面上绘制矩形草图
    swModel.Extension.SelectByID2 "Front Plane", _
        "PLANE", 0, 0, 0, False, 0, Nothing, 0
    swModel.SketchManager.InsertSketch True
    swModel.SketchManager.CreateCornerRectangle _
        -width / 2, -height / 2, 0, _
         width / 2,  height / 2, 0
    swModel.SketchManager.InsertSketch True

    ' 拉伸为实体
    swModel.FeatureManager.FeatureExtrusion2 _
        True, False, False, 0, 0, _
        depth / 1000, 0, _
        False, False, False, False, _
        0, 0, False, False, False, False, _
        True, True, True, 0, 0, False
End Sub
\end{lstlisting}

\begin{compare}[title={VBA 与文本脚本引擎的关键差异}]
\begin{tabularx}{\textwidth}{lXX}
\toprule
\textbf{维度} & \textbf{VBA (COM)} & \textbf{Lua/Python/Tcl (C API)} \\
\midrule
绑定机制 & COM \Tp{IDispatch}（后期绑定） & C 函数指针 / 栈操作 \\
类型系统 & COM VARIANT（运行时类型） & 脚本语言原生类型 \\
对象模型 & 宿主定义 COM 对象层次 & 宿主手动注册函数和类型 \\
跨平台 & 仅 Windows（COM 依赖） & 跨平台 \\
安全模型 & 宏安全策略（信任中心） & 沙箱 VM（Lua/AS） \\
性能 & COM 调用开销大（$\mu$s 级/调用） & 栈操作快（ns 级/调用） \\
用户群 & 办公人员、CAD 工程师 & 游戏开发者、系统程序员 \\
\bottomrule
\end{tabularx}
\end{compare}

\begin{warning}[title={VBA 的安全隐患}]
VBA 宏是恶意软件传播的常见载体（宏病毒）。Microsoft 已逐步收紧 VBA 宏的安全策略：Office 365 默认阻止来自互联网的宏执行（Mark-of-the-Web），AutoCAD 也引入了 \Tp{SECURELOAD} 机制限制自动加载的 LISP/VBA 脚本。在新项目中选择嵌入方案时，应评估是否需要更安全的沙箱替代（如 Office Add-in 使用 JavaScript/Web 技术）。
\end{warning}


\section{其他值得关注的引擎}\label{sec:other-engines}

\begin{center}
\begin{tabularx}{\textwidth}{lXXl}
\toprule
\textbf{引擎} & \textbf{语言} & \textbf{典型应用} & \textbf{体积} \\
\midrule
Duktape & JavaScript (ES5) & 嵌入式设备、游戏 & $\sim$250\,KB \\
QuickJS & JavaScript (ES2023) & Node.js 替代、嵌入式 & $\sim$500\,KB \\
Wren & Wren (Smalltalk-like) & 游戏（TIC-80 引擎） & $\sim$200\,KB \\
MiniScript & MiniScript & 教育（Mini Micro） & $\sim$300\,KB \\
MRuby & Ruby (子集) & 嵌入式 Web 服务 & $\sim$1\,MB \\
\bottomrule
\end{tabularx}
\end{center}


% ======================================================
\part{集成模式与架构}
% ======================================================

% ====================================
\chapter{函数绑定与调用约定}\label{ch:binding}
% ====================================

\section{绑定层设计模式}\label{sec:binding-patterns}

所有脚本集成方案的核心挑战都是 \textbf{绑定层}——在 C++ 的静态类型世界与脚本语言的动态类型世界之间架桥。主要有三种设计模式：

\subsection{手动栈操作}\label{sec:manual-stack}

直接使用脚本引擎的 C API（Lua C API、CPython API、Squirrel API）手动管理栈和类型转换。代码量大、容易出错，但性能最优：

\cpplabel
\begin{lstlisting}[style=cppstyle]
// 手动栈操作示例（Lua）
static int l_greet(lua_State* L) {
    const char* name = luaL_checkstring(L, 1);
    int age = luaL_checkinteger(L, 2);

    char buf[128];
    snprintf(buf, sizeof(buf),
             "Hello %s, age %d!", name, age);

    lua_pushstring(L, buf);
    return 1;
}
\end{lstlisting}

\subsection{模板元编程自动绑定}\label{sec:template-binding}

sol2 和 pybind11 使用 C++ 模板元编程自动推导参数类型和返回类型，将手动栈操作封装为声明式接口：

\cpplabel
\begin{lstlisting}[style=cppstyle]
// sol2 自动绑定：编译器自动生成栈操作代码
lua["greet"] = [](const std::string& name, int age) {
    return fmt::format("Hello {}, age {}!", name, age);
};
// 编译器推导：string → lua_tostring, int → lua_tointeger
// 返回值：string → lua_pushstring
\end{lstlisting}

\subsection{IDL/代码生成}\label{sec:idl-binding}

对于大规模绑定，可以使用接口描述语言（IDL）或代码生成工具自动产出绑定代码。典型工具包括 SWIG（支持 Python/Lua/多种目标语言）和 tolua++。

\section{调用约定对比}\label{sec:call-conv}

不同脚本引擎的函数调用机制差异显著：

\begin{center}
\begin{tabularx}{\textwidth}{lXX}
\toprule
\textbf{引擎} & \textbf{参数传递} & \textbf{返回值} \\
\midrule
Lua & 虚拟栈（索引取值） & 压栈 + 返回计数 \\
Python & \Tp{PyObject*} 引用 + \Tp{PyTuple} & \Tp{PyObject*} 新引用 \\
AngelScript & 预分配内存 + 指针 & 预分配输出参数 \\
ChaiScript & \Tp{Boxed\_Value} 类型擦除 & \Tp{Boxed\_Value} \\
Squirrel & 虚拟栈（类似 Lua） & 压栈 + SQRESULT \\
\bottomrule
\end{tabularx}
\end{center}

\section{类型映射策略}\label{sec:type-mapping}

\begin{guideline}[title={通用类型映射原则}]
\begin{enumerate}
  \item \textbf{基本类型}（int、float、bool、string）：所有引擎都支持直接映射，注意整数精度（Lua 5.3 区分 integer 和 float）。
  \item \textbf{容器类型}（vector、map、set）：pybind11 和 sol2 通过 \Tp{stl.h} 头文件自动转换；手动绑定需要逐元素遍历。
  \item \textbf{自定义类}：通过 userdata（Lua）、\Tp{PyObject}（Python）或 value type（AngelScript）映射。需要处理构造、析构和 GC 的协调。
  \item \textbf{智能指针}：sol2 和 pybind11 原生支持 \Tp{shared\_ptr} 的共享所有权语义。手动绑定时需要自行实现引用计数协调。
  \item \textbf{回调函数}：脚本函数可以映射为 \Tp{std::function}，但注意脚本 VM 的生命周期必须长于回调的使用期。
\end{enumerate}
\end{guideline}


% ====================================
\chapter{热重载与状态管理}\label{ch:hot-reload}
% ====================================

\section{热重载机制}\label{sec:hot-reload}

热重载是脚本集成最核心的价值——在不重启宿主程序的情况下更新脚本逻辑。实现方式因引擎而异：

\subsection{Lua 热重载}\label{sec:lua-hot-reload}

Lua 的热重载最直接：重新 \Fn{luaL\_dofile} 即可覆盖全局函数。保留状态需要额外处理：

\cpplabel
\begin{lstlisting}[style=cppstyle]
class ScriptManager {
    lua_State* L;
    std::filesystem::file_time_type last_modified;
    std::string script_path;

public:
    ScriptManager(const std::string& path)
        : script_path(path) {
        L = luaL_newstate();
        luaL_openlibs(L);
        reload();
    }

    bool reload() {
        auto mtime = std::filesystem::
            last_write_time(script_path);
        if (mtime <= last_modified) return false;

        // 保存旧状态
        lua_getglobal(L, "__saved_state");

        // 重新加载脚本
        if (luaL_dofile(L, script_path.c_str())
                != LUA_OK) {
            fprintf(stderr, "reload error: %s\n",
                    lua_tostring(L, -1));
            lua_pop(L, 2);
            return false;
        }

        // 恢复状态到新环境
        if (!lua_isnil(L, -1)) {
            lua_setglobal(L, "__saved_state");
        } else {
            lua_pop(L, 1);
        }

        last_modified = mtime;
        return true;
    }

    ~ScriptManager() { lua_close(L); }
};
\end{lstlisting}

\subsection{Python 模块重载}\label{sec:python-hot-reload}

Python 提供 \Fn{importlib.reload} 但限制较多——已有实例不会自动更新类定义：

\cpplabel
\begin{lstlisting}[style=cppstyle]
// C++ 侧触发 Python 模块重载
py::exec(R"(
    import importlib
    import game_logic
    importlib.reload(game_logic)
    # 注意：已有的 game_logic 实例不会自动更新
    # 需要手动重新初始化
)");
\end{lstlisting}

\subsection{AngelScript 模块重建}\label{sec:as-hot-reload}

AngelScript 支持丢弃并重建整个模块，是最干净的热重载方案：

\cpplabel
\begin{lstlisting}[style=cppstyle]
// 丢弃旧模块，重新编译
asIScriptModule* mod = engine->GetModule(
    "game", asGM_ALWAYS_CREATE); (*@\textcolor{gray}{// 丢弃旧模块}@*)

// 重新添加脚本节并编译
mod->AddScriptSection("main", new_source_code);
if (mod->Build() >= 0) {
    // 编译成功，新模块就绪
    // 旧模块中导出的对象将被 GC 回收
} else {
    // 编译失败，模块处于空状态
    // 可以加载备份字节码恢复
}
\end{lstlisting}

\section{状态持久化策略}\label{sec:state-persist}

热重载的核心挑战不是重新加载代码，而是 \textbf{保留运行时状态}。常用策略：

\begin{bestpractice}[title={状态持久化方案}]
\begin{enumerate}
  \item \textbf{全局状态表}：将所有可变状态存储在一个专用的全局表中（Lua）或字典中（Python），重载后自动保留。
  \item \textbf{序列化快照}：重载前将关键状态序列化为 JSON/二进制，重载后反序列化恢复。
  \item \textbf{事件驱动架构}：脚本只负责事件处理逻辑，状态由 C++ 侧管理。重载脚本不影响状态。
  \item \textbf{版本化状态}：为状态添加版本号，重载时执行状态迁移脚本（类似数据库 schema 迁移）。
\end{enumerate}
\end{bestpractice}

\section{文件监控}\label{sec:file-watch}

实现自动热重载需要监控脚本文件变化。跨平台方案：

\cpplabel
\begin{lstlisting}[style=cppstyle]
#include <filesystem>
#include <thread>
#include <chrono>

namespace fs = std::filesystem;

// 简单的轮询式文件监控（适合开发阶段）
class FileWatcher {
    fs::path path_;
    fs::file_time_type last_write_;
    std::function<void()> on_change_;

public:
    FileWatcher(fs::path p, std::function<void()> cb)
        : path_(std::move(p))
        , last_write_(fs::last_write_time(path_))
        , on_change_(std::move(cb)) {}

    void poll() {
        auto current = fs::last_write_time(path_);
        if (current != last_write_) {
            last_write_ = current;
            on_change_();
        }
    }
};
// 在主循环中每帧调用 watcher.poll()
\end{lstlisting}


% ======================================================
\part{C\# 脚本化方案}
% ======================================================

% ====================================
\chapter{Roslyn Scripting}\label{ch:roslyn}
% ====================================

C\# 的脚本化能力主要来自 Roslyn 编译器平台。
\seqsplit{Microsoft.CodeAnalysis.CSharp.Scripting} 包允许在运行时编译和执行 C\# 代码片段，相当于在 C\# 程序中内嵌了一个 C\# 解释器。

\section{基本用法}\label{sec:roslyn-basic}

\csharplabel
\begin{lstlisting}[style=csharpstyle]
using Microsoft.CodeAnalysis.CSharp.Scripting;
using Microsoft.CodeAnalysis.Scripting;

// 执行简单的 C# 表达式
int result = await CSharpScript.EvaluateAsync<int>(
    "1 + 2 * 3");
// result = 7

// 执行多行代码
var state = await CSharpScript.RunAsync(@"
    var items = new List<string> { ""a"", ""b"", ""c"" };
    var count = items.Count;
");
int count = state.GetVariable<int>("count");
// count = 3
\end{lstlisting}

\section{宿主对象绑定}\label{sec:roslyn-host}

Roslyn Scripting 允许向脚本注入宿主对象（Globals），脚本可以访问宿主对象的属性和方法：

\csharplabel
\begin{lstlisting}[style=csharpstyle]
// 定义宿主对象类型
public class GameGlobals
{
    public Player Player { get; set; }
    public IGameWorld World { get; set; }
    public void Log(string message) =>
        Console.WriteLine($"[Script] {message}");
}

public class Player
{
    public string Name { get; set; } = "Hero";
    public int HP { get; set; } = 100;
    public void TakeDamage(int amount) =>
        HP = Math.Max(0, HP - amount);
}

// 执行脚本，注入宿主对象
var globals = new GameGlobals {
    Player = new Player(),
    World = gameWorld
};

var options = ScriptOptions.Default
    .AddReferences(typeof(Player).Assembly)
    .AddImports("System", "System.Linq");

await CSharpScript.RunAsync(@"
    Player.TakeDamage(30);
    Log($""{Player.Name} HP: {Player.HP}"");
    // 输出: [Script] Hero HP: 70
", options, globals);
\end{lstlisting}

\section{编译缓存}\label{sec:roslyn-cache}

Roslyn 脚本首次执行需要编译（\textasciitilde 100--500\,ms），后续可以通过缓存编译结果避免重复编译：

\csharplabel
\begin{lstlisting}[style=csharpstyle]
// 预编译脚本为委托（首次约 200ms）
var script = CSharpScript.Create<int>(@"
    var sum = 0;
    for (int i = 0; i < 100; i++) sum += i;
    return sum;
");

// 编译一次
var compiled = script.CreateDelegate();

// 后续调用极快（微秒级）
int result1 = compiled();  // 4950
int result2 = compiled();  // 4950
\end{lstlisting}


% ====================================
\chapter{反射与动态类型}\label{ch:reflection-dynamic}
% ====================================

\section{反射}\label{sec:reflection}

C\# 的反射系统允许在运行时发现和调用类型成员，无需编译时绑定：

\csharplabel
\begin{lstlisting}[style=csharpstyle]
using System.Reflection;

// 运行时加载程序集
var asm = Assembly.LoadFrom("plugins/calculator.dll");
var calcType = asm.GetType("Calculator.Arithmetic");

// 动态创建实例
object instance = Activator.CreateInstance(calcType);

// 发现并调用方法
var addMethod = calcType.GetMethod("Add",
    new[] { typeof(double), typeof(double) });
double result = (double)addMethod.Invoke(
    instance, new object[] { 3.14, 2.86 });
// result = 6.0

// 枚举所有公开方法
foreach (var method in calcType.GetMethods(
    BindingFlags.Public | BindingFlags.Instance))
{
    Console.WriteLine($"  {method.ReturnType.Name} " +
        $"{method.Name}({string.Join(", ",
            method.GetParameters()
                .Select(p => p.ParameterType.Name))})");
}
\end{lstlisting}

\section{dynamic 关键字}\label{sec:dynamic}

C\# 的 \Tp{dynamic} 关键字推迟类型检查到运行时，简化了与动态语言的交互和插件系统：

\csharplabel
\begin{lstlisting}[style=csharpstyle]
// dynamic 跳过编译时类型检查
dynamic obj = GetScriptResult();
string name = obj.Name;     // 运行时解析
int result = obj.Calculate(10, 20); // 运行时解析

// 配合 ExpandoObject 实现动态对象
dynamic config = new ExpandoObject();
config.Host = "localhost";
config.Port = 8080;
config.Timeout = TimeSpan.FromSeconds(30);

// 运行时添加方法
((IDictionary<string, object>)config)["Validate"]
    = new Func<bool>(() => config.Port > 0);

if (config.Validate())
    Console.WriteLine($"Connecting to " +
        $"{config.Host}:{config.Port}");
\end{lstlisting}

\section{Source Generator}\label{sec:source-gen}

C\# 的 Source Generator（源生成器）在编译时自动生成代码，可以替代部分脚本化的功能：

\csharplabel
\begin{lstlisting}[style=csharpstyle]
// 源生成器在编译时自动产出代码
[Generator]
public class CommandGenerator : ISourceGenerator
{
    public void Execute(GeneratorExecutionContext ctx)
    {
        // 扫描标记了 [Command] 特性的方法
        // 自动生成注册和路由代码
        var source = @"
namespace Generated {
    public static class CommandRegistry {
        public static void RegisterAll(
            ICommandBus bus) {
            bus.Map(""attack"",
                (IGameState s, int dmg) =>
                    s.Player.TakeDamage(dmg));
            bus.Map(""heal"",
                (IGameState s, int amt) =>
                    s.Player.Heal(amt));
        }
    }
}";
        ctx.AddSource("CommandRegistry.g.cs", source);
    }
}
\end{lstlisting}

\begin{compare}[title={C\# 脚本化能力 vs C++ 脚本集成}]
C\# 的 Roslyn Scripting 允许在运行时编译和执行 \textbf{与宿主同一种语言} 的代码，这是 C++ 无法做到的（C++ 没有标准化的运行时编译器）。C\# 的反射、\Tp{dynamic}、表达式树和 Source Generator 共同构成了一个完整的动态执行体系，而 C++ 必须借助外部脚本语言来实现类似功能。
\end{compare}


% ======================================================
\part{综合对比与总结}
% ======================================================

% ====================================
\chapter{C++ 与 C\# 脚本化对比}\label{ch:cpp-vs-csharp}
% ====================================

\section{架构差异}\label{sec:arch-diff}

C++ 和 C\# 在脚本化方面存在根本性的架构差异：

\begin{compare}[title={根本差异}]
\begin{tabularx}{\textwidth}{lXX}
\toprule
\textbf{维度} & \textbf{C++} & \textbf{C\#} \\
\midrule
脚本语言 & 外部语言（Lua/Python 等） & 自身（Roslyn） \\
类型系统桥接 & 跨语言映射（静态 $\leftrightarrow$ 动态） & 同语言（可选 dynamic） \\
编译时开销 & 脚本零编译，宿主不变 & Roslyn 首次编译 $\sim$200\,ms \\
运行时依赖 & 额外嵌入 VM（Lua/Python） & .NET Runtime（已有） \\
工具链 & 各自独立（Lua 工具 + C++ 工具） & 统一（IDE、调试器、分析器） \\
调试 & 需跨语言调试器或日志 & 统一调试体验 \\
\bottomrule
\end{tabularx}
\end{compare}

\section{减少编译的效果}\label{sec:reduce-compile}

两种方案在减少编译方面的效果对比：

\begin{center}
\begin{tabularx}{\textwidth}{lXX}
\toprule
\textbf{场景} & \textbf{C++ + Lua/Python} & \textbf{C\# + Roslyn} \\
\midrule
修改一行逻辑 & 保存脚本文件 $\to$ 立即生效 & 保存 .csx 文件 $\to$ Roslyn 编译 $\to$ 生效 \\
首次加载开销 & Lua: $<$1\,ms; Python: $\sim$100\,ms & Roslyn: $\sim$200\,ms \\
增量更新 & 文件级重载 & 重新编译脚本片段 \\
热重载可靠性 & 需要手动处理状态迁移 & 同样需要处理状态迁移 \\
性能影响 & 脚本执行慢于原生 C++ & Roslyn 编译后的 IL 接近原生性能 \\
\bottomrule
\end{tabularx}
\end{center}

\section{适用场景分析}\label{sec:use-cases}

\begin{bestpractice}[title={选型建议}]
\begin{itemize}
  \item \textbf{游戏开发（热更新、AI、关卡逻辑）}：C++ + Lua（sol2），LuaJIT 性能足够，体积小，行业标准。
  \item \textbf{UE 项目开发（UE5）}：蓝图处理 UI、关卡逻辑和快速原型；C++ 处理核心系统。UE6+ 项目应优先采用 Verse 作为用户侧脚本，蓝图逐步退居可视化前端。
  \item \textbf{工具/编辑器扩展}：C++ + Python（pybind11），利用 Python 生态（NumPy、Pillow 等）。
  \item \textbf{EDA/CAD 工业工具}：C/C++ + Tcl，与现有工具链兼容，大量脚本资产可复用。
  \item \textbf{Office 办公自动化}：VBA 宏用于 Excel 报表生成、Word 批量处理、Access 数据操作；新项目可考虑 Office Add-in（JavaScript）作为跨平台替代。
  \item \textbf{CAD 参数化设计}：AutoCAD/SolidWorks/CATIA 中用 VBA 或 AutoLISP 实现批量绘图、参数化建模和设计自动化。
  \item \textbf{实时系统（嵌入式、仿真）}：C++ + AngelScript，静态类型减少运行时错误，字节码预编译加速启动。
  \item \textbf{快速原型/DLC 内容}：C++ + ChaiScript，语法接近 C++ 降低切换成本。
  \item \textbf{嵌入式 Web/IoT}：Duktape 或 QuickJS（嵌入式 JavaScript），适合需要 JavaScript 生态但资源受限的场景。
  \item \textbf{学术研究/DSL 设计}：Guile/Scheme，宏系统允许在语言层面定义新语法，适合编译器研究和领域特定语言。
  \item \textbf{C\#/.NET 项目}：优先使用 Roslyn Scripting（同语言，零学习成本）；若需要沙箱隔离则考虑 IL 动态生成。
  \item \textbf{插件系统}：C\# 用反射 + \Tp{Assembly.LoadFrom}；C++ 用脚本引擎或动态库（\Fn{dlopen}/\Fn{LoadLibrary}）。
\end{itemize}
\end{bestpractice}


% ====================================
\chapter{总结}\label{ch:conclusion}
% ====================================

脚本语言集成是 C/C++ 项目减少编译开销、提升开发迭代速度的成熟方案。核心结论：

\begin{enumerate}
  \item \textbf{Lua + sol2 是最通用的选择}：体积小、性能好（LuaJIT）、绑定简洁，是游戏行业的事实标准。
  \item \textbf{UE 生态正经历脚本范式转换}：UE6+ 将 Verse 定位为用户侧脚本主力，蓝图退居可视化前端——新项目应提前布局 Verse 技能栈。
  \item \textbf{Python + pybind11 适合需要丰富生态的场景}：科学计算、数据分析、机器学习相关的 C++ 项目首选。
  \item \textbf{VBA 仍是 Office/CAD 自动化的标准方案}：尽管安全隐患促使 Microsoft 收紧宏策略，VBA 在存量系统和工业工具链中的地位短期不可替代。
  \item \textbf{工业领域由 Tcl 主导}：EDA/CAD 工具链积累了大量 Tcl 脚本资产，新项目如需集成现有工具链应首选 Tcl。
  \item \textbf{AngelScript 适合对类型安全要求高的场景}：静态类型在编译时捕获错误，语法对 C++ 开发者友好。
  \item \textbf{C\# 拥有内在优势}：Roslyn Scripting 允许运行时编译执行同语言代码，反射和 dynamic 提供完整的动态能力——这是 C++ 无法匹敌的语言平台优势。
  \item \textbf{热重载的关键不是脚本引擎，而是状态管理}：无论选择哪种方案，都需要从一开始就设计好状态持久化和迁移策略。
\end{enumerate}


\appendix

% ====================================
\chapter{C++20 字面量类型与脚本注册优化}\label{app:literal-type}
% ====================================

\section{问题背景}\label{sec:lt-problem}

传统的脚本引擎绑定层在注册 C++ 类型和函数时，普遍使用运行时字符串作为标识符。例如 sol2 中调用 \seqsplit{lua.new\_usertype<T>("Player")}，或 pybind11 中调用 \seqsplit{py::class\_<Player>(m, "Player")}——这些字符串在运行时进行哈希查找和匹配。对于大型项目（数百个类型和数千个函数），注册表和查找表带来不可忽视的运行时开销。

C++20 引入了 \textbf{结构化类型的字面量模板参数}（Structural Type NTTP），允许将自定义类型用作非类型模板参数（Non-Type Template Parameter）。利用这一特性，可以将字符串编码为类型的一部分，使不同的注册名称在编译时就成为不同的类型——消除运行时字符串比较。

\section{TString 字面量类型}\label{sec:tstring}

核心工具是一个固定长度的字符串字面量类型：

\cpplabel
\begin{lstlisting}[style=cppstyle]
#include <algorithm>
#include <cstddef>
#include <string_view>

template <std::size_t N>
struct TString {
    char str[N]{};

    constexpr TString(const char (&s)[N]) {
        std::copy(s, s + N, str);
    }

    constexpr operator std::string_view() const {
        return {str, N - 1}; (*@\textcolor{gray}{// 排除末尾 '\textbackslash 0'}@*)
    }

    constexpr bool operator==(const TString&) const
        = default;
};

// 推导指引：从字符串字面量推导 TString<N>
template <std::size_t N>
TString(const char (&)[N]) -> TString<N>;
\end{lstlisting}

\Tp{TString} 满足 C++20 结构化类型（Structural Type）的要求：所有成员公开、无用户自定义析构、支持 \Tp{operator==}。因此它可以作为模板的非类型参数。

\section{类型级字符串区分}\label{sec:type-discrim}

利用 \Tp{TString} 作为模板参数，同一个函数模板对不同字符串实参产生不同的类型实例：

\cpplabel
\begin{lstlisting}[style=cppstyle]
// 通用注册器：Name 是编译期字符串，T 是目标类型
template <TString Name, typename T>
struct TypeRegistration {
    static constexpr auto name = Name;
    using type = T;

    // 每个 Name+T 组合生成唯一的 static 成员
    // 编译器保证不同类型不会共享地址
    static inline const int registration_id =
        next_id();

private:
    static int next_id() {
        static int counter = 0;
        return ++counter;
    }
};

// 不同字符串 → 不同类型 → 不同注册
using Reg1 = TypeRegistration<"Player", Player>;
using Reg2 = TypeRegistration<"Enemy",  Enemy>;
using Reg3 = TypeRegistration<"Item",   Item>;

// Reg1, Reg2, Reg3 是完全不同的类型
// 即使 T 相同，只要 Name 不同就是不同类型
static_assert(!std::is_same_v<
    TypeRegistration<"Vec2", Vec2>,
    TypeRegistration<"Vec3", Vec2>>);
\end{lstlisting}

\section{编译期注册表}\label{sec:compile-reg}

结合 \Tp{TString} NTTP 和模板特化，可以在编译期构建类型安全的注册表：

\cpplabel
\begin{lstlisting}[style=cppstyle]
// 编译期类型查找（通过字符串名 → 类型）
template <TString Name>
struct TypeRegistry;  (*@\textcolor{gray}{// 主模板未定义}@*)

// 每个类型显式注册（编译期特化）
template <>
struct TypeRegistry<"Player"> {
    using type = Player;
    static constexpr auto methods =
        std::tuple{&Player::GetHP,
                   &Player::TakeDamage};
};

template <>
struct TypeRegistry<"Enemy"> {
    using type = Enemy;
    static constexpr auto methods =
        std::tuple{&Enemy::AI\_Update,
                   &Enemy::GetTarget};
};

// 编译期查找（查找失败 → 编译错误，而非运行时崩溃）
template <TString Name>
using Lookup = typename TypeRegistry<Name>::type;

// 使用
static_assert(std::is_same_v<
    Lookup<"Player">, Player>);
\end{lstlisting}

\section{对脚本嵌入的优化效果}\label{sec:lt-benefit}

\begin{guideline}[title={字面量类型优化脚本嵌入}]
\begin{enumerate}
  \item \textbf{消除运行时字符串哈希}：注册表从 \Tp{unordered\_map<string, TypeInfo>} 变为编译期模板特化，查找从 $O(1)$ 哈希降为 $O(1)$ 编译期实例化。
  \item \textbf{类型安全保证}：错误的类型名称在编译时报错（模板未定义），而非运行时返回空指针。
  \item \textbf{零开销绑定代码生成}：结合 C++20 的 \Tp{consteval} 函数，可以在编译期生成完整的绑定胶水代码（类型转换、参数打包），运行时零额外开销。
  \item \textbf{减少二进制体积}：编译期注册表不需要存储字符串到类型的映射数据，减少只读数据段（\Tp{.rodata}）的大小。
\end{enumerate}
\end{guideline}

\cpplabel
\begin{lstlisting}[style=cppstyle]
// consteval 编译期绑定代码生成
template <TString Name, typename T>
consteval auto generate_bindings() {
    using Reg = TypeRegistry<Name>;
    // 编译期生成：类型名、方法数量、签名信息
    struct BindingInfo {
        std::string_view name;
        std::size_t method_count;
        bool has_constructor;
    };
    return BindingInfo{
        .name = Name,
        .method_count = std::tuple_size_v<
            decltype(Reg::methods)>,
        .has_constructor = std::is_default_
            constructible_v<T>
    };
}

// 编译期生成，运行时常量
constexpr auto player_info =
    generate_bindings<"Player", Player>();
// player_info.name == "Player"
// player_info.method_count == 2
// player_info.has_constructor == true
\end{lstlisting}

\begin{note}[title={适用边界}]
C++20 字面量类型 NTTP 适用于 \textbf{注册数量在编译期确定} 的场景。如果脚本接口需要在运行时动态注册（如插件系统、用户自定义类型），仍须使用传统的运行时哈希表。此外，\Tp{TString} 的长度是模板参数的一部分——\Tp{TString<6>} 和 \Tp{TString<7>} 是不同类型，这在某些场景下需要注意。
\end{note}


% ====================================
\chapter{C/C++ 作为嵌入语言}\label{app:cpp-embedded}
% ====================================

\section{工业仿真与仪器控制中的 C/C++ 嵌入}\label{sec:cpp-industrial}

在多数人的认知中，"脚本嵌入"意味着在 C/C++ 宿主中嵌入 Lua、Python 等高级语言。但在工业仿真和仪器控制领域，存在一个相反的模式：\textbf{C/C++ 本身被嵌入到更高层次的宿主系统中}，作为用户可扩展的计算后端。

\begin{compare}[title={C/C++ 作为嵌入语言的典型案例}]
\begin{tabularx}{\textwidth}{lX}
\toprule
\textbf{宿主系统} & \textbf{C/C++ 嵌入方式} \\
\midrule
Ansys (Mechanical/Fluent/HFSS) & 用户通过 \seqsplit{User-Defined Functions} (UDF) 或 \seqsplit{User Programmable Features} (UPF) 编写 C/C++ 代码，编译为共享库后加载到求解器中 \\
LabVIEW & Code Interface Node (CIN) 允许在图形化数据流图中嵌入 C 代码节点，编译为动态库后由 LabVIEW 运行时调用 \\
MATLAB/Simulink & MEX 文件（C/C++ 编写的扩展模块），S-Function（Simulink 中的 C/C++ 自定义模块） \\
COMSOL Multiphysics & 通过 External Material Model 接口嵌入 C++ 材料模型代码 \\
ABAQUS & \seqsplit{User Subroutines} 用 Fortran/C 编写自定义本构关系、接触算法等 \\
\bottomrule
\end{tabularx}
\end{compare}

\section{嵌入机制}\label{sec:cpp-embed-mechanism}

\subsection{Ansys UDF 示例}\label{sec:ansys-udf}

\cpplabel
\begin{lstlisting}[style=cppstyle]
// Ansys Fluent UDF: 自定义边界条件
#include "udf.h"

// 定义非均匀入口速度分布
DEFINE_PROFILE(inlet_velocity, thread, position) {
    face_t f;
    real y;   (*@\textcolor{gray}{// Ansys 内置 real 类型}@*)
    real x[ND_ND]; (*@\textcolor{gray}{// 坐标数组}@*)

    begin_f_loop(f, thread) {
        F_CENTROID(x, f, thread); (*@\textcolor{gray}{// 获取面心坐标}@*)
        y = x[1];
        // 抛物线速度分布: v = vmax * (1 - (y/H)^2)
        F_PROFILE(f, thread, position) =
            10.0 * (1.0 - pow(y / 0.05, 2.0));
    }
    end_f_loop(f, thread)
}

// 编译: 在 Fluent TUI 中执行
//   Define > User-Defined > Functions > Compiled
//   添加源文件 → Build → Load
\end{lstlisting}

\subsection{LabVIEW CIN 机制}\label{sec:labview-cin}

LabVIEW 的 CIN（Code Interface Node）允许用户在图形化数据流图中嵌入 C 函数。用户编写符合 LabVIEW CIN 接口的 C 代码，编译为平台特定的动态库后，LabVIEW 在运行时通过固定入口点调用：

\cpplabel
\begin{lstlisting}[style=cppstyle]
// LabVIEW CIN: 自定义信号处理函数
// CIN 接口规范（由 LabVIEW 生成头文件）
#include "extcode.h" (*@\textcolor{gray}{// LabVIEW 提供的头文件}@*)

// CIN 函数签名必须严格匹配
// LabVIEW 中定义的输入/输出类型
CIN MgErr CINRun(
    float64* inputArray,
    int32    inputSize,
    float64* outputArray,
    int32*   outputSize)
{
    // 自定义 DSP 算法（例如：快速傅里叶变换）
    // 这部分用 C 而非 LabVIEW 图形化编程
    // 因为 FFT 的位反转和蝶形运算用文本
    // 代码表达远比图形化节点直观

    *outputSize = inputSize;
    for (int32 i = 0; i < inputSize; i++) {
        outputArray[i] = /* FFT 计算 */ ;
    }
    return noErr;
}
\end{lstlisting}

\section{图形化编程的局限性}\label{sec:gpl-limits}

LabVIEW 本身是图形化编程（G 语言）的代表。虽然图形化编程在简单数据流和仪器控制场景中直观易用，但当逻辑复杂度上升时，其劣势显著暴露：

\begin{warning}[title={图形化编程的复杂度瓶颈}]
\begin{itemize}
  \item \textbf{可读性急剧下降}：当数据流连线超过数百条时，VI（Virtual Instrument）的图形变得如同"意大利面"——交叉的连线比文本代码的缩进更难追踪。一个 50 行的 C 函数等价于可能铺满 3--4 屏的图形化节点。
  \item \textbf{版本控制困难}：图形化代码以二进制格式存储（\Tp{.vi} 文件），传统的 diff/merge 工具无法有效工作，团队协作和代码审查成本极高。
  \item \textbf{算法表达力受限}：递归、高阶函数、泛型编程等文本语言中的基础抽象在图形化环境中要么不支持、要么极其笨拙。这就是为什么 LabVIEW 提供了 CIN 机制——承认图形化编程在复杂算法场景中的不足。
  \item \textbf{调试体验}：图形化调试依赖动画高亮数据流动，对于并发逻辑和竞态条件的排查效率远低于文本调试器的断点和调用栈查看。
\end{itemize}
\end{warning}

\section{与脚本嵌入方案的对比}\label{sec:cpp-vs-script}

\begin{compare}[title={C/C++ 嵌入 vs 脚本语言嵌入}]
\begin{tabularx}{\textwidth}{lXX}
\toprule
\textbf{维度} & \textbf{C/C++ 嵌入} & \textbf{Lua/Python/VBA 嵌入} \\
\midrule
目标用户 & 高级开发/算法工程师 & 设计师/分析师/普通用户 \\
编译需求 & 需要 C/C++ 编译工具链 & 无需编译，直接解释执行 \\
热更新 & 需要重新编译动态库并重新加载 & 保存脚本文件即可生效 \\
性能 & 接近原生（因为就是原生） & 解释执行，慢 5--100$\times$ \\
开发门槛 & 高（需要 C/C++ 和宿主 API 知识） & 低（语法简洁，API 文档友好） \\
调试 & 需要附加调试器到宿主进程 & 脚本级调试器或 print 调试 \\
安全风险 & 段错误可能崩溃宿主进程 & 沙箱 VM 隔离脚本错误 \\
适用场景 & 高性能计算核心、自定义物理模型 & 业务逻辑、UI、报表、配置 \\
\bottomrule
\end{tabularx}
\end{compare}

\begin{bestpractice}[title={工程实践建议}]
C/C++ 嵌入方案的合理性建立在两个前提上：(1) 性能要求确实需要原生代码（如 CFD 求解器中的自定义湍流模型、有限元分析中的非线性材料本构）；(2) 使用者具备高级 C/C++ 开发能力。对于不满足这两个条件的场景，Lua、Python、VBA 或 Blueprint 等脚本方案是更优选择——它们将开发效率提升一到两个数量级，而性能损失在多数业务逻辑中可以忽略。

对于 Ansys/LabVIEW 等工业软件的用户，推荐的层次化策略是：用内置脚本语言（APDL、G 语言）处理常规自动化和流程控制；用 C/C++ 嵌入（UDF、CIN）实现性能关键的自定义算法；两者互补而非替代。
\end{bestpractice}


\end{document}
